\documentclass[11pt,a4paper]{article}
\usepackage[T1]{fontenc}
\usepackage[utf8]{inputenc}
\usepackage{lmodern,microtype}
\usepackage[margin=25mm]{geometry}
\usepackage{amsmath,amssymb,amsthm,mathtools,bm}
\usepackage{booktabs,tabularx,array,longtable,enumitem}
\usepackage{xcolor}
\usepackage[colorlinks=true,linkcolor=blue!50!black,citecolor=blue!50!black,urlcolor=blue!50!black]{hyperref}
\usepackage[nameinlink,noabbrev]{cleveref}
\numberwithin{equation}{section}
\newtheorem{theorem}{Theorem}[section]
\newtheorem{proposition}[theorem]{Proposition}
\newtheorem{lemma}[theorem]{Lemma}
\newtheorem{corollary}[theorem]{Corollary}
\theoremstyle{definition}
\newtheorem{definition}[theorem]{Definition}
\newtheorem{assumption}[theorem]{Assumption}
\theoremstyle{remark}
\newtheorem{remark}[theorem]{Remark}
\DeclareMathOperator{\Id}{Id}
\DeclareMathOperator{\Var}{Var}
\DeclareMathOperator{\rank}{rank}
\DeclareMathOperator{\spanop}{span}
\DeclareMathOperator{\diag}{diag}
\DeclareMathOperator{\atanh}{atanh}
\DeclareMathOperator{\range}{range}
\newcommand{\R}{\mathbb R}
\newcommand{\E}{\mathbb E}
\newcommand{\Lie}{\mathcal L}
\newcommand{\Dop}{\mathcal D}
\newcommand{\norm}[1]{\left\lVert #1\right\rVert}
\newcommand{\abs}[1]{\left\lvert #1\right\rvert}
\newcommand{\pair}[2]{\left\langle #1,#2\right\rangle}
\DeclareUrlCommand\file{\urlstyle{tt}}
\setlength{\emergencystretch}{2em}
\setlist{itemsep=3pt,topsep=5pt}

\title{Intrinsic Finite-Degree Identifiability of Shallow\\Randomized-Signature Reservoirs}
\author{}
\date{13 September 2026}

\begin{document}
\maketitle
\begin{abstract}
For shallow biased componentwise-tanh vector fields with $d$ inputs and
$D$ state coordinates, we construct an explicit finite-parameter witness
whose degree-$L$ Hall evaluation matrix at zero has rank
$\min(D,\nu_{d,L})$. Here $\nu_{d,L}$ is the dimension of the free step-$L$
Lie algebra. A classical grading-derived affine action and finite weighted
jets give a block upper triangular tanh Hall matrix with diagonal blocks
$\tfrac12 I$ and exact selected minors at every width. The established
analytic-minor method then gives the same rank outside a Lebesgue-null
analytic exceptional set in the full weight and bias space.
Consequently, one zero-initialized vector of the
degree-$L$ truncated response locally determines all Hall coordinates
when $D\geq\nu_{d,L}$; the width is sharp for continuous local recovery
from this observation model. This is a statement about a finite
polynomial, not an exact inverse of a full CDE endpoint.
Separate supporting results give a scalar two-channel homogeneous
certificate through degree four, its rank-seven cross-degree obstruction,
ideal scale-jet reconstruction, and an explicit eight-observation scalar
truncated inverse with rational stability constants. A finite remainder
bound extends that scalar decoder to approximate recovery from eight
actual CDE responses. A larger-width tensor-to-tanh construction proves
independence of all finite-degree word-response fields. These statements
have distinct observation assumptions and quantitative content; they
do not imply a linear-readout theorem or explain the preliminary
financial experiments that motivated the question.
\end{abstract}

\tableofcontents

\section{Introduction}
A shallow randomized-signature reservoir maps an input path into a state
vector using fixed nonlinear vector fields. Which intrinsic path
coordinates can that vector retain? We answer this question for the
degree-$L$ truncated response, whose unknowns are the
$\nu_{d,L}$ Hall coordinates of a truncated log-signature. The observation
is one $D$-dimensional vector at a fixed nonzero input scale, with
initialization fixed at zero.

Generic bracket generation for this shallow architecture has already been
studied by Monzio Compagnoni et al.\ \cite{MonzioCompagnoni}.
Cuchiero, Larsson, and Teichmann \cite{CuchieroLarssonTeichmann} use
analytic minors of Lie-field evaluations to prove generic interpolation,
and Grong and Vega-Molino \cite{GrongVegaMolino} use the scalar constant/tanh
pair for landmark controllability. Here we specify the Hall family through
a prescribed finite degree and construct exact maximal minors for it.
The distinction is the finite-degree realization and observation theorem,
rather than bracket generation or analytic genericity by themselves.

Our main result is the generic rank law
\[
 \rank\bigl[\rho_\theta(H_1)(0)\ \cdots\
                 \rho_\theta(H_{\nu_{d,L}})(0)\bigr]
 =\min(D,\nu_{d,L})
\]
for shallow biased componentwise tanh in the full parameter space.
The Hall basis is ordered by degree, and $\rho_\theta$ replaces
generators by vector fields and brackets by vector-field brackets.
This matrix is the linearization of the truncated response, with
nonzero degree-dependent scale factors. It yields local analytic
recovery at width $D\geq\nu_{d,L}$. Invariance of domain rules out a
continuous locally injective encoding of an open Hall-coordinate
neighborhood into fewer than $\nu_{d,L}$ real outputs.

The proof uses a classical affine action obtained from an invertible
grading derivation. We realize its leading weighted terms by actual
shallow tanh fields at finite parameters. A finite-jet argument preserves
the diagonal Hall blocks at zero, while allowing nonlinear entries above
them. Graded quotients supply the smaller-width witnesses, and inactive
coordinates supply larger widths. The resulting nonzero minor is real
analytic in all weights and biases. This proves genericity at the fixed
state directly, without sampling initial states or taking a limiting
parameter as the witness. The witness-and-analytic-minor method has a
direct precedent in \cite[Lemma 6.1]{CuchieroLarssonTeichmann}; the
construction here supplies the prescribed-degree minor for fixed tanh.

\paragraph{Contributions and observation models.}
The central contribution is the explicit witness and sharp fixed-state
Hall-rank law in \cref{thm:hallrank}, together with the one-vector local
truncated inverse in \cref{cor:hallinverse} and its dimensional lower
bound. Three supporting parts address different questions. First, the
previously used constant/tanh pair has homogeneous injectivity through
degree four but joint function rank seven, demonstrating why homogeneous
information alone need not give fixed-state recovery. Second, ideal
normalized scale derivatives give a general graded decoder, while a
designed set of eight scalar truncated observations has an exact local
stability certificate; a uniform remainder and a clipped decoder give
approximate recovery from the corresponding eight actual CDE responses.
Third, a tensor-to-tanh transfer at the larger sufficient width
$\sum_{m=0}^L d^m$ proves independence of all nonempty word responses
through length $L$, a larger family than the Hall fields.

We distinguish ideal scale jets (Object A), finite truncated response
polynomials (Object B), and actual CDE responses (Object C). The new
one-vector inverse concerns Object B. It inherits neither the scalar
design's constants nor its actual-response estimate. It also does not
turn a nonlinear inverse into a trained linear readout.

Preliminary pricing and partial-information hedging experiments
motivated the representation question: useful history features do not
by themselves establish preservation of specified path coordinates.
The existing financial material is retained in \cref{sec:finance} as
motivation, without a performance implication from the Lie theorems.
The classical affine construction and the bounded comparison with
reconstruction and randomized-signature sources are discussed in
\cref{sec:related}.

\begin{table}[htbp]
\centering
\caption{Notation. The symbol $N$ is reserved for a time-step count.}
\label{tab:notation}
\begin{tabularx}{\textwidth}{@{}lX@{}}
\toprule
Symbol & Meaning \\
\midrule
$d,D,L,N$ & Input channels, state dimension, truncation level, time steps \\
$m,r,h,z$ & Homogeneous degree, input scale, common scale multiplier, Hall coordinates \\
$T(E),T^{\leq L}(E)$ & Tensor algebra and its degree-$L$ truncation, $E=\R^d$ \\
$\Lie_{d,m},n_m,\nu_{d,L}$ & Homogeneous free-Lie space, its dimension, $\sum_{m=1}^L n_m$ \\
$H_{m,a},H_a$ & Homogeneous Hall basis and its flattened degree order \\
$H_\theta$ & Matrix of Hall-field evaluations at the fixed state zero \\
$\mathcal E,\Lambda_a$ & Grading derivation and $\mathcal E^{-1}\operatorname{ad}_a\mathcal E$ \\
$\Gamma,\rho$ & Tensor-to-operator homomorphism and Lie-to-vector-field map \\
$F_w=\Gamma(w)\Id$ & Word-response field, distinct from the operator $\Gamma(w)$ \\
$S_m,\ell_m$ & Signature and log-signature components of degree $m$ \\
$J_m$ & Normalized scale jet $(1/m!)\partial_r^mY_T^{y,r}|_{r=0}$ \\
$F_y^{\leq L}(r;z)$ & Degree-$L$ truncated response polynomial \\
$B,K,\Delta(R)$ & Scalar-design Jacobian, $\norm{B^{-1}}_\infty$, variation majorant \\
$D_0,Q_{d,L}$ & $\sum_{m=0}^L d^m$ and $\sum_{m=1}^L d^m$ \\
\bottomrule
\end{tabularx}
\end{table}

\section{Algebraic and analytic preliminaries}
\label{sec:prelim}
\subsection{Words, signatures, and Hall coordinates}
Fix integers $d,L\geq1$ and put $E=\R^d$, with ordered basis letters
$X_1,\ldots,X_d$. A word $w=i_1\cdots i_m$ denotes
$X_{i_1}\otimes\cdots\otimes X_{i_m}$ and has length $|w|=m$.
The empty word is the unit $1$. The tensor algebra
$T(E)=\bigoplus_{m\geq0}E^{\otimes m}$ has concatenation product $uv$.
Its truncation $T^{\leq L}(E)$ discards products of degree greater than $L$;
$\pi_m$ is projection onto degree $m$.

For a continuous bounded-variation (BV) path $X:[0,T]\to\R^d$, define
\begin{equation}\label{eq:signature}
 S_m(X)=\sum_{i_1,\ldots,i_m}
 \left(\int_{0<t_1<\cdots<t_m<T}
 dX_{t_1}^{i_1}\cdots dX_{t_m}^{i_m}\right)i_1\cdots i_m,
 \qquad S_0(X)=1.
\end{equation}
The path and the formal letters share the conventional symbol $X$; their
roles are distinguished by time subscripts and algebraic expressions.
The earliest time is on the left. For concatenated paths $X*\widehat X$,
splitting each integration simplex at the joining time yields Chen's
identity $S^{\leq L}(X*\widehat X)=S^{\leq L}(X)S^{\leq L}(\widehat X)$.
For a straight segment with increment $a\in E$, the simplex volume is
$1/m!$, so its signature is $\sum_{m=0}^L a^m/m!$.

The free Lie algebra $\Lie_d$ is realized as the Lie subalgebra of $T(E)$
generated by the letters under $[u,v]=uv-vu$. Its homogeneous components are
$\Lie_{d,m}=\Lie_d\cap E^{\otimes m}$ and
$\Lie_d^{\leq L}=\bigoplus_{m=1}^L\Lie_{d,m}$. Here the latter notation
first denotes a vector space. Equipping it with a truncated bracket makes
the free step-$L$ nilpotent Lie algebra, but the nonlinear vector-field map
below need not respect that truncated bracket.

A nonempty word is Lyndon if it is lexicographically smaller than all its
nontrivial cyclic rotations. A letter brackets to itself. A longer Lyndon
word $w=uv$ brackets as $H_w=[H_u,H_v]$, with $v$ its longest proper Lyndon
suffix. Order these words by degree and then lexicographically. The
free-Lie embedding theorem and the Lyndon basis theorem state that this
construction gives a basis in each degree; its size is
\begin{equation}\label{eq:witt}
 n_m=\dim\Lie_{d,m}=\frac1m\sum_{k\mid m}\mu(k)d^{m/k},
 \qquad \nu_{d,L}=\sum_{m=1}^L n_m,
\end{equation}
where $\mu$ is the arithmetic M\"obius function. These are classical inputs
\cite[Theorems 0.4--0.5, 4.9, Corollary 4.14, Theorem 5.1]{Reutenauer}.
They apply to our finite ordered alphabet over $\R$. No algebraic freeness
or basis theorem is inferred from a numerical rank.

Write the ordered degree-$m$ basis as $H_{m,1},\ldots,H_{m,n_m}$. A coordinate
vector $z=(z_{m,a})\in\R^{\nu_{d,L}}$ represents
$\ell(z)=\sum_{m=1}^L\ell_m(z)$, where
$\ell_m(z)=\sum_a z_{m,a}H_{m,a}$. In the truncated algebra define
\[
 \exp_\otimes(\ell)=\sum_{k=0}^L\frac{\ell^k}{k!},\qquad
 \log_\otimes(1+s)=\sum_{k=1}^L\frac{(-1)^{k+1}s^k}{k}.
\]
Terms of degree above $L$ are discarded. These inverse finite polynomial
maps are used only for positive-degree $\ell,s$.

The precise signature realization input needed later is as follows: the
step-$L$ signature of a continuous BV path belongs to
$\exp_\otimes(\Lie_d^{\leq L})$, and every element of that group is the
signature of a piecewise-linear path. In the Friz--Victoir author draft,
Theorems 7.28 and 7.30 identify this group with products of exponentials of
degree-one elements \cite{FrizVictoir}. Each such product is realized by
concatenating straight segments, in the order of \eqref{eq:signature}.
Consequently $\ell(X)=\log_\otimes S^{\leq L}(X)$ has the coordinates just
defined. This realization statement concerns unrestricted BV inputs.

\subsection{Operators and represented Lie fields}
Let $U\subseteq\R^D$ be open and let $V_1,\ldots,V_d$ be smooth vector
fields on $U$. For a scalar or vector-valued smooth function $f$, put
$\Dop_Vf=Df\,V$, component by component. Define
\begin{equation}\label{eq:Gamma}
 \Gamma(i_1\cdots i_m)=\Dop_{V_{i_1}}\cdots\Dop_{V_{i_m}},
 \qquad \Gamma(1)=\mathrm{identity\ operator}.
\end{equation}
Composition acts on the function to its right; the first letter is the
outermost operator. Thus $\Gamma(uv)=\Gamma(u)\Gamma(v)$ on the full
tensor algebra.

For two vector fields the convention is
$[V,W]=DW\,V-DV\,W$. Direct differentiation gives
\[
 (\Dop_V\Dop_W-\Dop_W\Dop_V)f
 =Df(DW\,V-DV\,W)=\Dop_{[V,W]}f,
\]
since the two Hessian terms cancel. The universal property of the free Lie
algebra therefore defines $\rho:\Lie_d\to\mathcal X(U)$ by
$\rho(X_i)=V_i$, where $\mathcal X(U)$ is the space of smooth fields with
this bracket. Induction on brackets proves the operator identity
\begin{equation}\label{eq:Lieoperator}
 \Gamma(H)=\Dop_{\rho(H)}\quad(H\in\Lie_d),\qquad
 \Gamma(H)\Id=\rho(H).
\end{equation}
The second equality follows from the first, not conversely. We call
$F_w=\Gamma(w)\Id$ a \emph{word-response field}; it is a vector-valued
function, whereas $\Gamma(w)$ is a differential operator. A word-response
field generally is not itself a represented Lie field.

\begin{definition}[Three independence requirements]
Homogeneous free-Lie injectivity through degree $L$ means that each
$\rho_m=\rho|_{\Lie_{d,m}}$ is injective. Injectivity on the truncated
free-Lie vector space means that $\rho|_{\Lie_d^{\leq L}}$ is injective
as a linear map. Finite-degree word-response field independence means that
the family $\{\Gamma(w)\Id:1\leq|w|\leq L\}$ is jointly linearly
independent over $\R$ as functions on the observation domain.
\end{definition}
By the tensor embedding and \eqref{eq:Lieoperator}, the last property implies
the second, which implies the first. None of these finite statements says
that represented brackets of degree greater than $L$ vanish.

\subsection{Finite evaluations and analytic arguments}
\begin{lemma}[Scalar point-coordinate evaluation]\label{lem:evaluation}
Let $f_1,\ldots,f_s:A\to\R^D$ be linearly independent functions on a
nonempty set $A$. There exist pairs $(y_j,k_j)\in A\times\{1,\ldots,D\}$,
$1\leq j\leq s$, such that the matrix
$(f_a(y_j)_{k_j})_{j,a}$ is nonsingular.
\end{lemma}
\begin{proof}
Let $W$ be their $s$-dimensional span. Scalar point-coordinate evaluations
are linear functionals on $W$. Their common kernel is zero: a function whose
every coordinate vanishes everywhere is the zero function. Their span in
$W^*$ must therefore have dimension $s$, since otherwise its annihilator
in $W$ would be nonzero. Choose $s$ evaluations forming a basis of $W^*$.
Their matrix on the given basis of $W$ is nonsingular.
\end{proof}

\begin{lemma}[Analytic determinant certificate]\label{lem:analytic}
Let $\Theta$ be a connected open subset of a finite-dimensional Euclidean
space, and let $f:\Theta\to\R$ be real analytic with $f(\theta_0)\ne0$
at some point. Then $\{f\ne0\}$ is open, dense, and has full Lebesgue
measure in $\Theta$.
\end{lemma}
\begin{proof}
Continuity proves openness. At a zero of a nontrivial analytic function,
some finite derivative is nonzero, by the analytic identity principle on a
connected domain. A derivative of one lower order vanishes there and has
nonzero gradient. The zero set is thus covered by the countably many
regular zero sets of derivatives. Each is locally a smooth hypersurface by
the implicit function theorem and is Lebesgue null. A countable chart cover
gives nullity. A null set contains no open ball, proving density. This is
the analytic zero-set theorem in the form of Mityagin \cite{Mityagin}.
\end{proof}

\begin{lemma}[Local analytic inverse]\label{lem:inverse}
Let $F$ be real analytic near $0\in\R^s$ with nonsingular derivative
$B=DF(0)$. There are neighborhoods of $0$ and $F(0)$ on which $F$ is a
bijection with real-analytic inverse.
\end{lemma}
\begin{proof}
Translate the image and multiply by $B^{-1}$ to write the map as
$z+R(z)$ with $R(0)=DR(0)=0$. Its convergent real power series extends
holomorphically to a small complex ball. Shrink that ball until
$\sup\norm{DR}\leq1/2$. For sufficiently small $w$, the iteration
$z_{j+1}=w-R(z_j)$ maps the closed ball into itself and is a contraction.
It has a unique fixed point, which solves $z+R(z)=w$. The iterates are
holomorphic in $w$ and converge uniformly on a smaller ball. The uniform
limit theorem for holomorphic functions gives a holomorphic inverse there.
Restricting to real $w$ gives the stated inverse; uniqueness also gives
local injectivity. The neighborhood can depend on every coefficient of $F$.
\end{proof}

\subsection{The truncated observation and its linearization}
\begin{definition}[Truncated response: Object B]
For $z\in\R^{\nu_{d,L}}$, $y\in U$, and $r\in\R$, define
\begin{equation}\label{eq:truncated}
 F_y^{\leq L}(r;z)=y+\sum_{m=1}^L r^m
 \Gamma\bigl(\pi_m\exp_\otimes\ell(z)\bigr)\Id(y).
\end{equation}
This is a finite polynomial in $z$ and $r$. A scalar observation selects
one coordinate at a chosen state and scale.
\end{definition}
All products in this definition are first expanded and truncated by
tensor word length, equivalently by Hall weighted degree. The resulting
finite representative lies in the full tensor algebra, where $\Gamma$
is defined; the nonlinear fields need not act on the truncated quotient.
The empty word has known coefficient one and contributes $\Id(y)=y$.
At zero it contributes zero, so it is not an additional unknown
coordinate in any evaluation-rank claim.

Differentiating \eqref{eq:truncated} at $z=0$ eliminates all products of
two or more Lie factors. For a basis element $H_{m,a}$ its derivative is
\begin{equation}\label{eq:zero_derivative}
 \partial_{z_{m,a}}F_y^{\leq L}(r;0)_k
 =r^m\rho(H_{m,a})(y)_k.
\end{equation}
This identity uses the coefficient one of the linear term in the tensor
exponential. It introduces no additional factorial. The ideal normalized
scale jets and the actual CDE response will be defined separately in
\cref{sec:graded}.

\section{Sharp Hall rank at initialization zero}
\label{sec:hallrank}
Fix integers $d,D,L\geq1$ and consider the full architecture
\begin{equation}\label{eq:architecture}
 V_i^\theta(x)=\tanh(A_ix+b_i),\qquad
 \theta=((A_i,b_i))_{i=1}^d\in
 \mathcal P_{d,D}:=\R^{d(D^2+D)},\quad x\in\R^D.
\end{equation}
Here $A_i\in\R^{D\times D}$, $b_i\in\R^D$, and tanh acts componentwise.
Write $\nu=\nu_{d,L}$ and flatten the existing degree-ordered basis as
$H_1,\ldots,H_\nu$, with $m_a=\deg H_a$. Let
$\mathfrak g=\bigoplus_{m=1}^L\mathfrak g_m$ be the free step-$L$ Lie
algebra. For the witness construction we write its generators as $e_i$,
the images of the formal letters $X_i$. When evaluated on nonlinear
fields, $H_a$ always means its canonical representative in the full
free Lie algebra.

\subsection{The grading-derived affine action}
\begin{lemma}[Affine action from the grading]\label{lem:gradingaffine}
Let $\mathfrak a$ be a finite positively graded real Lie algebra and
$\mathcal E|_{\mathfrak a_m}=mI$. For $a\in\mathfrak a$ and $c\in\R$ set
\[
 \Lambda_a=\mathcal E^{-1}\operatorname{ad}_a\mathcal E,\qquad
 W_a^c(x)=ca-\Lambda_ax.
\]
Here $\operatorname{ad}_a b=[a,b]$, and matrix brackets use
$[\Lambda_a,\Lambda_b]=\Lambda_a\Lambda_b-\Lambda_b\Lambda_a$.
Then $[\Lambda_a,\Lambda_b]=\Lambda_{[a,b]}$,
$\Lambda_ab-\Lambda_ba=[a,b]$, and, with the convention of
\eqref{eq:Lieoperator}, $[W_a^c,W_b^c]=W_{[a,b]}^c$.
\end{lemma}
\begin{proof}
Positive grading makes $\mathcal E$ an invertible derivation. Conjugating
the adjoint commutator proves the first identity. The derivation identity
gives
\[
 \Lambda_ab-\Lambda_ba
 =\mathcal E^{-1}\bigl([a,\mathcal Eb]+[\mathcal Ea,b]\bigr)=[a,b].
\]
As $DW_a^c=-\Lambda_a$, direct substitution gives
\[
 [W_a^c,W_b^c](x)
 =c(\Lambda_ab-\Lambda_ba)
   +(\Lambda_b\Lambda_a-\Lambda_a\Lambda_b)x
 =c[a,b]-\Lambda_{[a,b]}x.
\]
Thus a Hall element $H$ acts as $W_H^c$ and evaluates to $cH$ at zero.
For homogeneous $a,b$ of degrees $s,t$, respectively,
$\Lambda_ab=t[a,b]/(s+t)$.
\end{proof}
This is the classical nonsingular-cocycle construction with adjoint
module and cocycle $\mathcal E$; see Burde
\cite[Section 3.1, Proposition 3.3, Corollaries 3.4--3.5]{BurdePreLie}.
The calculation above fixes the vector-field sign used here.

\subsection{A finite-parameter tanh witness}
\begin{lemma}[Leading Hall blocks of the tanh witness]\label{lem:hallwitness}
Identify $\mathfrak g$ with $\R^\nu$ in its ordered Hall basis. With
$c=1/2$ and $q=\atanh(1/2)$, the admissible fields
\begin{equation}\label{eq:hallwitness}
 V_i(x)=\tanh(qe_i-\Lambda_{e_i}x)
\end{equation}
satisfy, for $H\in\mathfrak g_m$ and grade projection $P_k$,
\[
 P_k\rho_V(H)(0)=0\quad(k>m),\qquad
 P_m\rho_V(H)(0)=cH.
\]
Their degree-ordered Hall matrix is block upper triangular with diagonal
blocks $cI_{n_m}$ and determinant $2^{-\nu}$.
For a positively graded quotient $\overline{\mathfrak g}$ with a retained
Hall basis whose generator images are coordinate vectors or zero, the
same leading-block identities hold on the retained columns, with
determinant $2^{-\dim\overline{\mathfrak g}}$.
\end{lemma}
\begin{proof}
Give each grade-$k$ coordinate weight $k$. The map $\Lambda_{e_i}$
raises grade by one. Hence the grade-one field components are the
constants $ce_i$. Each grade-$k$ component for $k\geq2$ is
$\tanh(\ell_{i,a}(x))$, where
$\ell_{i,a}=-(\Lambda_{e_i}x)_a$ is linear in grade-$(k-1)$ coordinates.
Its linear term agrees with $W_{e_i}^c$.

Only finite jets are needed. With $F_w=\Gamma(w)\Id$, the recursion
$F_{iv}=DF_v\,V_i$ shows that $F_w(0)$ for $|w|=m$ uses ordinary
generator derivatives through order $m-1$. Replace the fields by their
Taylor polynomials through order $L$ at zero. All evaluations through
length $L$ are unchanged. Apart from the affine terms, these polynomials
contain only $\alpha_n\ell_{i,a}^{\,n}$ with odd $3\leq n\leq L$.
No assertion about an infinite-radius Taylor series is needed.

A monomial operator $p(x)\partial_{x_a}$ has weight shift $s-k$ when
$p$ has weighted degree $s$ and $x_a$ has weight $k$; shifts add under
composition. Every affine operator term has shift $-1$. Each nonlinear
term has shift
\[
 n(k-1)-k=-1+(n-1)(k-1)\geq1.
\]
A length-$m$ word acting on a grade-$k$ coordinate therefore has an
all-affine contribution of weight $k-m$; any contribution containing a
nonlinear term has strictly larger weight. Evaluation at zero selects
weight zero. The value vanishes if $k>m$ and agrees with the affine
value if $k=m$. Expand a homogeneous Hall element into its length-$m$
words and apply \cref{lem:gradingaffine} and \eqref{eq:Lieoperator}.
This gives $cH$ in grade $m$, and the asserted determinant follows.
The argument is unchanged in the stated graded quotient. Components in
grades below $m$ can be nonzero; the full matrix need not equal $cI$.
\end{proof}

\subsection{The rank law in every state dimension}
\begin{theorem}[Generic sharp fixed-state Hall rank]\label{thm:hallrank}
For every choice of integers $d,D,L\geq1$, there is a nontrivial real-analytic function
$f_{d,D,L}$ on the full parameter space $\mathcal P_{d,D}$ in
\eqref{eq:architecture} such that
\begin{equation}\label{eq:hallrank}
 H_\theta=\bigl[\rho_\theta(H_1)(0)\ \cdots\
                    \rho_\theta(H_\nu)(0)\bigr],\qquad
 \rank H_\theta=\min(D,\nu)
 \quad\text{if }f_{d,D,L}(\theta)\ne0.
\end{equation}
The function is a fixed maximal-size minor, and an explicit finite
witness gives $f_{d,D,L}=2^{-\min(D,\nu)}$. Its zero set is Lebesgue
null. The rank law therefore holds almost surely under any jointly
absolutely continuous law on all weights and biases.
\end{theorem}
\begin{proof}
For $D=\nu$, use \eqref{eq:hallwitness} and \cref{lem:hallwitness}.
For $D<\nu$, let $s_p=\sum_{m=1}^p n_m$, with $s_0=0$, choose the
least $p$ with $s_p\geq D$, and put $t=D-s_{p-1}$. Retain all Hall
vectors below grade $p$ and the first $t$ in grade $p$. If $U_p$ spans
the remaining grade-$p$ vectors, then
\[
 I=U_p\oplus\bigoplus_{k=p+1}^L\mathfrak g_k
\]
is a graded ideal: bracketing a discarded vector with any positive
grade raises its degree above $p$, or gives zero. Thus
$\overline{\mathfrak g}=\mathfrak g/I$ has dimension $D$ and the
retained Hall images form its homogeneous basis.

Construct $\overline{\mathcal E}$ and
$\overline\Lambda_a=\overline{\mathcal E}^{-1}
\operatorname{ad}_a\overline{\mathcal E}$ on this quotient and use the
original $d$ generator labels in
\[
 V_i(x)=\tanh(q\overline e_i-
                    \overline\Lambda_{\overline e_i}x).
\]
The retained $D$ Hall columns and all rows have determinant $2^{-D}$
by \cref{lem:hallwitness}. If $D<d$, this construction retains the
first $D$ generators and kills the others; the retained fields are
constant $ce_i$ and the other fields are zero. Thus it also covers
that case. This argument constructs an ideal quotient before evaluating
fields, rather than arbitrarily deleting rows of a nonlinear witness.
Only the selected columns are needed: discarded Hall elements may have
nonzero nonlinear evaluations. The affine action respects the quotient;
no such factorization is asserted for the tanh action.

For $D>\nu$, embed \eqref{eq:hallwitness} in the first $\nu$ coordinates
and set the remaining weights and biases to zero. Word differentiation
preserves the inactive coordinates. The first $\nu$ rows and all Hall
columns retain determinant $2^{-\nu}$.

Fix these row and column selections and let $f_{d,D,L}$ be the resulting
minor for arbitrary $\theta$, always at initial state zero. The fields
are jointly real analytic in $(\theta,x)$; finite differentiation,
Hall expansion, evaluation at zero, and a determinant preserve
analyticity. The finite witnesses show that $f_{d,D,L}$ is nontrivial
on the connected Euclidean parameter space.
Lemma~\ref{lem:analytic} gives its open dense full-measure nonzero set.
There the minor supplies the lower rank bound, and the matrix dimensions
give the matching upper bound.
\end{proof}
The theorem does not automatically apply to bias-free, tied, sparse, or
other singular parameter laws. Its exceptional set is obtained at the fixed state
zero, not by taking a slice of a joint parameter--state genericity claim.

\subsection{One-vector local recovery and sharp width}
\begin{corollary}[Local inverse of the truncated response]\label{cor:hallinverse}
Fix integers $d,D,L\geq1$ with $D\geq\nu_{d,L}$, and a parameter with
$f_{d,D,L}(\theta)\ne0$ as in \cref{thm:hallrank}. For each fixed
$r\ne0$, the map $z\mapsto F_0^{\leq L}(r;z)$ from
\eqref{eq:truncated} permits local analytic recovery near $z=0$ of all
$\nu_{d,L}$ Hall coordinates from one zero-initialized vector.
For $D=\nu_{d,L}$ it has a local analytic inverse; for $D>\nu_{d,L}$
the selected $\nu_{d,L}$ output-coordinate projection has such an inverse.
\end{corollary}
\begin{proof}
Let $\delta_r\ell(z)=\sum_a r^{m_a}z_aH_a$. The observation is
\[
 F_0^{\leq L}(r;z)
 =\left(\Gamma_\theta\left(
   \pi_{\leq L}\sum_{j=0}^L
          \frac{(\delta_r\ell(z))^j}{j!}\right)\Id\right)(0).
\]
This is a finite tensor representative and its empty-word contribution
is zero. By \eqref{eq:zero_derivative},
\begin{equation}\label{eq:halljacobian}
 D_zF_0^{\leq L}(r;0)
 =H_\theta\diag(r^{m_1},\ldots,r^{m_\nu}).
\end{equation}
Let $P$ select the rows used in $f_{d,D,L}$, with $P=I$ for $D=\nu$.
The projected derivative has determinant
$r^{\sum_a m_a}f_{d,D,L}(\theta)\ne0$.
Lemma~\ref{lem:inverse} gives neighborhoods $U,\Omega$ of zero in
$\R^\nu$ and an analytic map $g:\Omega\to U$ with
$g(PF_0^{\leq L}(r;z))=z$ for $z\in U$.
For $D>\nu$ this inverts a projection, not a map between open sets of
unequal dimensions. The neighborhoods may depend on $\theta$ and $r$.
\end{proof}
The same nonzero-minor parameter set works for every $r\ne0$, but gives
no uniform inverse neighborhood or quantitative constant. The conclusion
concerns Object B; no exact inversion of an actual CDE endpoint or
linear-readout approximation follows.

\begin{proposition}[Lower bound on scalar outputs]\label{prop:dimension}
Let $O\subset\R^{\nu_{d,L}}$ be open and nonempty. If a continuous map
$F:O\to\R^s$ is injective on some nonempty open neighborhood, then
$s\geq\nu_{d,L}$.
\end{proposition}
\begin{proof}
The invariance-of-domain theorem states that a continuous injective map
from an open subset of $\R^n$ to $\R^n$ has open image
\cite[Theorem 2B.3]{Hatcher}. If $s<n=\nu_{d,L}$, restrict to the
neighborhood of injectivity and append $n-s$ zero coordinates to $F$.
This is a continuous injective map into $\R^n$. Its image lies in a
proper linear subspace, which has empty interior, contradicting that
theorem. All its hypotheses hold because the chosen domain is open and
the original map is continuous there.
\end{proof}

Applying \cref{prop:dimension} with $s=D$ and combining it with
\cref{cor:hallinverse} makes $D=\nu_{d,L}$ the sharp width for
continuous local recovery of an open Hall-coordinate neighborhood from
this one-vector truncated observation. This lower bound is not a
minimum number of runs for multiple-query schemes or a statement about
restricted path families.

\section{An exact scalar homogeneous theorem}
\label{sec:scalar}
Consider two scalar fields
\begin{equation}\label{eq:scalar}
 V_i(x)=\tanh(a_ix+b_i),\quad i=1,2,\qquad
 \theta=(a_1,b_1,a_2,b_2)\in\R^4.
\end{equation}
A scalar field is identified with its coefficient multiplying $\partial_x$.
Set
\begin{equation}\label{eq:witness}
 q=\frac{\log3}{2}=\atanh(1/2),\quad
 \theta_*=(1,0,0,q),\qquad V_1(x)=\tanh x,\quad V_2(x)=1/2.
\end{equation}
Up to exchanging generators and rescaling the constant field, this is
the pair in Grong and Vega-Molino \cite[Remark 3.2]{GrongVegaMolino}.
Their polynomial bracket argument gives landmark controllability.
Here we determine homogeneous free-Lie injectivity through degree four
and the exact mixed-degree relation for this pair; its use as a generator
pair is not a new construction.
The ordered Hall elements through degree four are
\begin{align}\label{eq:Hall8}
 H_1&=X_1,& H_2&=X_2,\nonumber\\
 H_3&=[X_1,X_2],&H_4&=[X_1,[X_1,X_2]],\nonumber\\
 H_5&=[[X_1,X_2],X_2],&H_6&=[X_1,[X_1,[X_1,X_2]]],\nonumber\\
 H_7&=[X_1,[[X_1,X_2],X_2]],&H_8&=[[[X_1,X_2],X_2],X_2].
\end{align}
Their Lyndon words are $1,2,12,112,122,1112,1122,1222$.
Their degrees are $(1,1,2,3,3,4,4,4)$. In \eqref{eq:Hall8} the aligned arrangement is
only typographical; the column order is always the numerical subscript.

\begin{theorem}[Scalar homogeneous injectivity through degree four]
\label{thm:scalar}
For the architecture \eqref{eq:scalar}, each map
$\rho_m:\Lie_{2,m}\to\mathcal X(\R)$, $1\leq m\leq4$, is injective
at \eqref{eq:witness}. There exists an open dense full-Lebesgue-measure
subset of $\R^4$ on which all four homogeneous maps are injective
separately. This does not assert injectivity on their direct sum.
\end{theorem}
\begin{proof}
Evaluate the degree-one fields at $(0,q)$, degree two at $0$, degree three
at $(-q,q)$, and degree four at $(-q,0,q)$. The exact matrices in the
ordering \eqref{eq:Hall8} are
\begin{align}\label{eq:Ms}
 M_1&=\begin{pmatrix}0&1/2\\1/2&1/2\end{pmatrix},&
 M_2&=\begin{pmatrix}-1/2\end{pmatrix},\nonumber\\
 M_3&=\begin{pmatrix}15/32&3/16\\15/32&-3/16\end{pmatrix},&
 M_4&=\begin{pmatrix}
 -57/128&-3/32&3/64\\-1/2&0&1/4\\-57/128&3/32&3/64
 \end{pmatrix}.
\end{align}
Their entries can be checked by the polynomial bracket recurrence in
\cref{sec:cross}. Direct determinant expansion gives
\begin{equation}\label{eq:dets}
 \begin{gathered}
 \det M_1=-\frac14,\quad\det M_2=-\frac12,\quad
 \det M_3=-\frac{45}{256},\quad\det M_4=\frac{135}{8192},\\[4pt]
 \prod_{m=1}^4\det M_m=-\frac{6075}{16777216}.
 \end{gathered}
\end{equation}
A linear relation among fields in one degree would vanish at its selected
points, so its coefficient vector would be in the kernel of a nonsingular
matrix. Thus each homogeneous map is injective at the witness.

For general $\theta$, use the same fixed initial points and the same Hall
columns to define $M_m(\theta)$. The map $(\theta,x)\mapsto V_i^\theta(x)$
is real analytic. Finite differentiation and the bracket formula preserve
joint real analyticity; hence the determinant product is analytic on
$\R^4$. Its value in \eqref{eq:dets} is nonzero. Apply
\cref{lem:analytic} to the product. On its nonzero locus all four matrices
are nonsingular, proving the stated simultaneous homogeneous conclusion.
\end{proof}

All entries in \eqref{eq:Ms} are rational because $\tanh(0)=0$ and
$\tanh(\pm q)=\pm1/2$. The points need not be rational to obtain a
rational certificate. Floating-point conditioning can aid witness discovery
but plays no role in this proof.

\section{The exact cross-degree obstruction}
\label{sec:cross}
Put $u=\tanh x$. If fields have scalar coefficients $P(u)$ and $Q(u)$,
then $du/dx=1-u^2$ and the bracket coefficient is
\begin{equation}\label{eq:polybracket}
 \mathcal B(P,Q)=(1-u^2)(P Q'-P'Q),
\end{equation}
where the primes on the right differentiate polynomials in $u$.
Applying this recurrence to $P_1=u$ and $P_2=1/2$ gives
\begin{align}\label{eq:P8}
 P_1&=u,&P_2&=\tfrac12,&P_3&=\tfrac12(u^2-1),\nonumber\\
 P_4&=\tfrac12(1-u^4),&P_5&=\tfrac12(u^3-u),\nonumber\\
 P_6&=\tfrac12(3u^6-3u^4+u^2-1),&
 P_7&=u^3-u^5,&P_8&=\tfrac14(1-4u^2+3u^4).
\end{align}
These satisfy $\rho_{\theta_*}(H_j)(x)=P_j(\tanh x)$.
For example, $\mathcal B(u,1/2)=-(1-u^2)/2$,
$\mathcal B(u,P_3)=(1-u^4)/2$, and
$\mathcal B(P_3,1/2)=(u^3-u)/2$. The remaining three brackets are
$\mathcal B(u,P_4)$, $\mathcal B(u,P_5)$, and
$\mathcal B(P_5,1/2)$, yielding precisely \eqref{eq:P8}. This gives all
eight fields and an exact degree bound of six.

\begin{proposition}[Joint rank and complete kernel]\label{prop:cross}
At $\theta_*$ the map $\rho:\Lie_2^{\leq4}\to\mathcal X(\R)$ has
rank seven. In the ordering \eqref{eq:Hall8} its kernel is
\begin{equation}\label{eq:kernel8}
 \ker\rho|_{\Lie_2^{\leq4}}=
 \spanop\{4H_3+3H_4+2H_8\}.
\end{equation}
\end{proposition}
\begin{proof}
The coefficient matrix of \eqref{eq:P8}, with rows indexed by
$1,u,\ldots,u^6$, is
\begin{equation}\label{eq:coeffmatrix}
 C=\begin{pmatrix}
 0&1/2&-1/2&1/2&0&-1/2&0&1/4\\
 1&0&0&0&-1/2&0&0&0\\
 0&0&1/2&0&0&1/2&0&-1\\
 0&0&0&0&1/2&0&1&0\\
 0&0&0&-1/2&0&-3/2&0&3/4\\
 0&0&0&0&0&0&-1&0\\
 0&0&0&0&0&3/2&0&0
 \end{pmatrix}.
\end{equation}
The determinant of its first seven columns is $-3/32$, so its rank is
seven. Multiplication by $(0,0,4,3,0,0,0,2)^\top$ gives zero. The kernel
is one dimensional by rank-nullity, proving completeness of the relation.
A polynomial vanishing after substitution $u=\tanh x$ vanishes on
$(-1,1)$ and is zero. Thus polynomial rank equals function rank.
\end{proof}

The relation uses degrees two, three, and four. It is consistent with the
four nonsingular homogeneous matrices. At a common nonzero input scale
$r$, the linearized Hall response fields are $r^{\deg H_j}P_j(\tanh x)$.
Their joint rank remains seven, since this multiplies columns by nonzero
constants. No number of initial-state evaluations at that one scale yields
a rank-eight zero-coordinate Jacobian. This obstructs a local analytic
inverse at zero with a nonsingular derivative. A singular derivative alone
does not prove failure of nonlinear local injectivity.

\section{Graded recovery from ideal scale jets}
\label{sec:graded}
\subsection{The observation object and its coefficient formula}
\begin{assumption}[Smooth fields and BV inputs]\label{ass:smooth}
The vector fields $V_1,\ldots,V_d$ are known and $C^\infty$ on a nonempty
open set $U\subseteq\R^D$. The input $X:[0,T]\to\R^d$ is continuous
and of bounded variation. Observations at an initial state $y\in U$ can
access input scaling in a neighborhood of zero.
\end{assumption}
For sufficiently small real $r$, consider
\begin{equation}\label{eq:CDE}
 dY_t^{y,r}=r\sum_{i=1}^dV_i(Y_t^{y,r})\,dX_t^i,
 \qquad Y_0^{y,r}=y.
\end{equation}
The scale changes the input to $rX$, or equivalently multiplies all fields
by $r$. It is not a rescaling of the solution state with unchanged fields.

\begin{definition}[Ideal normalized scale jets: Object A]
At an initial state $y$, the degree-$m$ ideal observation is
\begin{equation}\label{eq:jetdef}
 J_m(y;X)=\frac1{m!}\left.\partial_r^mY_T^{y,r}(X)\right|_{r=0}.
\end{equation}
One scalar jet observation selects one coordinate of this vector.
\end{definition}

\begin{proposition}[Finite Chen--Fliess coefficient identity]\label{prop:jets}
Under \cref{ass:smooth}, for each $y\in U$ and finite $L$, the solution
of \eqref{eq:CDE} exists for all sufficiently small $|r|$ through time $T$,
depends smoothly on $r$ near zero, and satisfies
\begin{equation}\label{eq:jets}
 Y_T^{y,r}=y+\sum_{m=1}^L r^m\Gamma(S_m(X))\Id(y)
 +O(|r|^{L+1}),\qquad
 J_m(y;X)=\Gamma(S_m(X))\Id(y).
\end{equation}
The solution neighborhood and remainder bound may depend on $y$, $X$, the
fields, and $L$.
\end{proposition}
\begin{proof}
Choose a closed ball $K=\overline B(y,a)\subset U$. On it all fields are
bounded by a common $M$ and Lipschitz with a common constant $K_1$ in a
fixed vector norm. Let $\omega=\sum_i\Var(X^i;[0,T])$. If $\omega=0$
the solution is constant and the assertion holds. Otherwise choose $|r|$
so that $|r|M\omega<a$ and $|r|K_1\omega<1$. The integral-equation map
\[
 Z\longmapsto y+r\sum_i\int_0^{\cdot}V_i(Z_t)\,dX_t^i
\]
maps the closed sup-norm ball of continuous paths centered at $y$ into
itself and is a contraction. Its fixed point remains in $K$. Shrinking
the scale interval gives a uniform contraction constant below one.
The integral map is smooth in the path and scale on a slightly smaller
ball: each derivative is the integral of a finite derivative of $V_i$
applied to continuous perturbations, bounded by a compact-state bound
times $\omega$. Difference quotients of the fixed-point equation converge
because its linearized operator has inverse given by the uniformly
convergent Neumann series. Repeating this argument for successive
derivatives, whose inhomogeneous terms involve only derivatives already
constructed, proves smooth dependence on $r$ to every fixed finite order.

For a smooth scalar function $f$, the BV chain rule gives
\[
 f(Y_t^{y,r})=f(y)+r\sum_i\int_0^t(\Dop_{V_i}f)(Y_s^{y,r})\,dX_s^i.
\]
Here the solution has bounded variation because its defining integral is
bounded by $|r|M\omega$. The chain rule follows by summing the first-order
Taylor formula along partitions: its error is bounded by a constant times
the largest oscillation of $Y$ on a partition interval times its total
variation, and hence tends to zero. This also verifies the chain rule's
continuity and finite-variation hypotheses in the present setting.
Substitute this identity successively for the functions
$\Dop_{V_i}f$. At length $m$, the earliest differential acts on the
outside, giving the operator ordering in \eqref{eq:Gamma}. After $L$
substitutions, the remainder is $r^{L+1}$ times a sum over words of length
$L+1$ of simplex integrals whose integrand is
$(\Gamma(i_1\cdots i_{L+1})f)(Y_{t_1}^{y,r})$. It is uniformly bounded
on $K$ for each of the finitely many words. If $C_f$ is their maximum
absolute bound, the absolute remainder is at most
\[
 |r|^{L+1}C_f\frac{\omega^{L+1}}{(L+1)!}.
\]
Indeed the sum of total-variation product measures over the ordered simplex
is $\omega^{L+1}/(L+1)!$; the variation measure is nonatomic because $X$
is continuous. Apply this calculation to each coordinate of $f=\Id$.
Smooth scale dependence and uniqueness of finite Taylor coefficients yield
\eqref{eq:jets}. The factor $1/m!$ in \eqref{eq:jetdef} is exactly the
Taylor normalization; the iterated integrals already contain their own
simplex factors.
\end{proof}

No convergence of an infinite Chen--Fliess series is needed. In particular,
the estimate proves a finite expansion at zero scale without making the
higher-degree terms available from a fixed nonzero-scale observation.

\subsection{Triangularity of tensor exponentiation}
For $\ell=\ell_1+\cdots+\ell_L\in\Lie_d^{\leq L}$ define
\begin{equation}\label{eq:Psidef}
 \Psi_m(y;\ell_1,\ldots,\ell_{m-1})=
 \sum_{k=2}^m\frac1{k!}
 \sum_{\substack{j_1+\cdots+j_k=m\\j_a\geq1}}
 \Gamma(\ell_{j_1}\cdots\ell_{j_k})\Id(y),
 \qquad\Psi_1=0.
\end{equation}
Every $j_a$ in this formula is at most $m-1$. We also define the formal
graded observation at $z$ by
$J_m^{\rm form}(y;z)=\Gamma(\pi_m\exp_\otimes\ell(z))\Id(y)$.

\begin{lemma}[Homogeneous tensor triangularity]\label{lem:triangular}
For $1\leq m\leq L$,
\begin{equation}\label{eq:expgraded}
 \pi_m\exp_\otimes(\ell)
 =\sum_{k=1}^m\frac1{k!}
 \sum_{\substack{j_1+\cdots+j_k=m\\j_a\geq1}}
 \ell_{j_1}\cdots\ell_{j_k}.
\end{equation}
Consequently
\begin{equation}\label{eq:triangular}
 J_m^{\rm form}(y;z)
 =\rho(\ell_m(z))(y)+\Psi_m(y;\ell_1(z),\ldots,\ell_{m-1}(z)).
\end{equation}
For actual paths, $J_m^{\rm form}(y;z(X))=J_m(y;X)$.
\end{lemma}
\begin{proof}
Expand each ordered product in $\ell^k=(\sum_j\ell_j)^k$. A term has
degree $j_1+\cdots+j_k$; positivity of all degrees makes $k\leq m$
necessary for a degree-$m$ term. This proves \eqref{eq:expgraded}, without
commuting its factors. The $k=1$ term is $\ell_m$. Apply $\Gamma$ and
$\Id$, use \eqref{eq:Lieoperator}, and collect $k\geq2$ terms into
\eqref{eq:Psidef}. The path assertion follows from
$S^{\leq L}(X)=\exp_\otimes\ell(X)$ and \cref{prop:jets}.
\end{proof}

For orientation, the positive-degree tensor components through four are
\begin{align}\label{eq:exp4}
 S_1&=\ell_1,\nonumber\\
 S_2&=\ell_2+\tfrac12\ell_1^2,\nonumber\\
 S_3&=\ell_3+\tfrac12(\ell_1\ell_2+\ell_2\ell_1)
       +\tfrac16\ell_1^3,\nonumber\\
 S_4&=\ell_4+\tfrac12(\ell_1\ell_3+\ell_2^2+\ell_3\ell_1)
 +\tfrac16(\ell_1^2\ell_2+\ell_1\ell_2\ell_1+\ell_2\ell_1^2)
 +\tfrac1{24}\ell_1^4.
\end{align}
These identities also specify the lower-order corrections in the scalar
decoder without choosing a second tensor convention.

\begin{assumption}[Homogeneous injectivity on the observation domain]
\label{ass:injective}
For every $1\leq m\leq L$, the map
$\rho_m:\Lie_{d,m}\to\mathcal X(U)$ is injective.
\end{assumption}
For general smooth fields this hypothesis concerns their restrictions to
$U$. Independence on a larger domain does not imply independence on every
smaller open set without an additional continuation property.

\begin{theorem}[Conditional graded Hall reconstruction]\label{thm:graded}
Under \cref{ass:injective}, there are $n_m$ point-coordinate pairs
$(y_{m,j},k_{m,j})$ in each nonzero degree $m$ such that the corresponding
scalar formal observations $J_m^{\rm form}(y_{m,j};z)_{k_{m,j}}$ recover
every $z\in\R^{\nu_{d,L}}$ recursively and exactly. Thus
$\nu_{d,L}$ scalar graded observations suffice. Under
\cref{ass:smooth} the same decoder recovers the truncated log-signature
of a BV path from its ideal scale jets.
\end{theorem}
\begin{proof}
The represented basis fields in degree $m$ are independent on $U$.
By \cref{lem:evaluation}, choose the stated pairs so that
\[
 E_m=\bigl(\rho(H_{m,a})(y_{m,j})_{k_{m,j}}\bigr)_{j,a=1}^{n_m}
\]
is nonsingular. Let $j_m$ collect the observed degree-$m$ scalars, and
let $\psi_m$ collect the same evaluations of \eqref{eq:Psidef}.
With $z_m=(z_{m,a})_a$, \eqref{eq:triangular} gives
\begin{equation}\label{eq:recursive}
 z_1=E_1^{-1}j_1,\qquad
 z_m=E_m^{-1}\bigl(j_m-\psi_m(z_1,\ldots,z_{m-1})\bigr).
\end{equation}
The right side uses only previously recovered blocks. Induction proves
existence of the exact decoder and uniqueness of its output for every
formal coordinate vector. When $n_m=0$ there is no block or observation
in that degree. Applying the actual-path coefficient identity proves the
last assertion. No smallness of the formal coordinates is required for
this algebraic recursion, although the definition of the actual jets uses
a scale neighborhood that may depend on the path.
\end{proof}

In the scalar case $E_m=M_m$. Their exact inverses and three rational
round-trip checks appear in \cref{app:graded}. The recursion separates the
mixed-degree relation in \eqref{eq:kernel8} by keeping its components in
different jet orders.

\begin{proposition}[A homogeneous kernel is formally invisible]
\label{prop:kernel}
Suppose a nonzero $H\in\Lie_{d,m}$, $m\leq L$, satisfies
$\rho(H)=0$ on $U$. For every $\ell\in\Lie_d^{\leq L}$ and $t\in\R$,
replacing $\ell$ by $\ell+tH$ leaves all formal graded observations
through $L$ unchanged, and leaves every degree-$L$ truncated response
unchanged. There are piecewise-linear paths with these respective
truncated log-signatures. Equality of their actual full CDE responses is
not asserted.
\end{proposition}
\begin{proof}
By \eqref{eq:Lieoperator}, $\Gamma(H)=\Dop_{\rho(H)}$ is the zero
operator on $U$. Expand the difference of the two tensor exponentials
into ordered products of homogeneous Lie factors. Every product contains
at least one factor $H$. Since $\Gamma$ is a homomorphism, its image
is a composition containing a zero operator and is zero on $U$.
The vanishing is on an open set, so differentiation by operators to its
left does not introduce a nonzero term. Homogeneity ensures that each
such product has a definite total degree: truncation retains or discards
the entire product, preserving the argument in every degree. Applying
the resulting operators to $\Id$ proves the observation assertions.
The realization input in \cref{sec:prelim} supplies piecewise-linear
paths for both exponentials. Their higher signature levels need not
coincide, so no conclusion about full responses follows.
\end{proof}
This is a necessity statement for the unrestricted formal graded model
and, with realization, for truncated path information. It does not supply
a continuous selection of realizing paths. It also does not apply to the
mixed-degree scalar relation by treating it as a homogeneous element.

\section{Finite state-scale observations of the truncated model}
\label{sec:finite}
The actual terminal solution $Y_T^{y,r}(X)$ is \emph{Object C}. For actual
path coordinates it agrees with \eqref{eq:truncated} up to the local
remainder in \eqref{eq:jets}. This estimate does not identify the
truncated polynomial from finitely many values of Object C. Object A
instead differentiates Object C ideally at zero scale, where the
coefficient identity is exact.

\begin{theorem}[Conditional finite truncated observation and inverse]
\label{thm:finite}
Assume \cref{ass:injective}. Let $I\subset\R\setminus\{0\}$ be any
nonempty open interval. There exist $\nu_{d,L}$ triples
$(r_j,y_j,k_j)\in I\times U\times\{1,\ldots,D\}$ such that the map
\[
 \Phi(z)=\bigl(F_{y_j}^{\leq L}(r_j;z)_{k_j}-(y_j)_{k_j}\bigr)
 _{j=1}^{\nu_{d,L}}:\R^{\nu_{d,L}}\longrightarrow\R^{\nu_{d,L}}
\]
has nonsingular derivative at zero and a local analytic inverse there.
The inverse concerns Object B and a neighborhood of zero coordinates.
\end{theorem}
\begin{proof}
By \eqref{eq:zero_derivative}, the derivative functions are
$r^m\rho(H_{m,a})(y)_k$.
Suppose a real linear combination of these functions vanishes on all
$(r,y,k)$. Grouping the coefficients by degree gives
$\sum_m r^m f_m(y)_k=0$. For fixed $y,k$ this is a polynomial vanishing
on an open interval, so every coefficient is zero. Thus $f_m=0$ on $U$
for every $m$. Homogeneous injectivity forces every original Hall
coefficient to vanish. The derivative functions are jointly independent.
Apply \cref{lem:evaluation}, or its scalar version on the set of triples,
to obtain a square nonsingular evaluation matrix. It is $D\Phi(0)$ by
\eqref{eq:zero_derivative}. The map $\Phi$ is polynomial in $z$ even
when the fields are only smooth in $y$, so \cref{lem:inverse} applies.
\end{proof}

The design can use several coordinates in a single run, or repeated
scales and states. The theorem only counts real outputs.
Thus \cref{thm:finite} matches an intrinsic scalar-output dimension bound.
It is not a universal minimum number of reservoir runs. It also does not
make the map globally invertible. The fixed-zero, one-vector conclusion
of \cref{cor:hallinverse} instead uses the stronger Hall evaluation rank
available in the stated architecture at $D\geq\nu_{d,L}$.

\section{An explicit eight-output scalar decoder}
\label{sec:eight}
Use the scalar witness \eqref{eq:witness} and the eight coordinates in
\eqref{eq:Hall8}. Let the ordered state-scale design be
\begin{equation}\label{eq:design}
 \mathcal D_*=\bigl((-q,1),(-q,2),(-q,3),(0,2),(0,3),
 (q,1),(q,2),(q,3)\bigr).
\end{equation}
For its pairs $(y_j,r_j)$ define the centered polynomial map
\begin{equation}\label{eq:Phi}
 \Phi_*(z)=\bigl(F_{y_j}^{\leq4}(r_j;z)-y_j\bigr)_{j=1}^8.
\end{equation}
The literal responses equal this vector plus the known initial-state
vector. Centering changes neither the derivative nor any reconstruction
difference. In particular it is not an approximation replacing $q$ by a
rational number.

\begin{theorem}[Fixed eight-output truncated inverse]\label{thm:eight}
At $\theta_*$, the zero-coordinate Jacobian $B=D\Phi_*(0)$ satisfies
\begin{equation}\label{eq:detB}
 \det B=-\frac{258733035}{524288}\ne0.
\end{equation}
Consequently \eqref{eq:Phi} has a local analytic inverse at zero. Keeping
the same ordered design, the corresponding truncated map for
\eqref{eq:scalar} has a local analytic inverse at zero on an open dense
full-Lebesgue-measure certificate set of parameters $\theta\in\R^4$.
\end{theorem}
\begin{proof}
By \eqref{eq:zero_derivative}, the $j,a$ entry of the actual derivative
of \eqref{eq:Phi} is $r_j^{\deg H_a}P_a(\tanh y_j)$. Substitution into
\eqref{eq:P8} gives the exact matrix in \cref{eq:Bfull}. Exact elimination,
or multiplication by the inverse displayed in \cref{app:inverse}, gives
\eqref{eq:detB}; all arithmetic is rational. This is a derivative of the
nonlinear truncated map, rather than an unrelated field-evaluation matrix.
\Cref{lem:inverse} proves the first inverse statement.

For arbitrary $\theta$, leave $y_j,r_j$ fixed and replace only the fields.
Finite differentiation of the jointly analytic fields makes every entry
of the zero-coordinate Jacobian analytic in $\theta$. Its determinant
is nonzero at $\theta_*$ by \eqref{eq:detB}. Apply
\cref{lem:analytic}, followed by \cref{lem:inverse} for each parameter
in the nonzero locus. No uniform inverse neighborhood is obtained.
\end{proof}

The intrinsic count is $\nu_{2,4}=2+1+2+3=8$. It explains why eight
scalar outputs are used. As a finite computational check of the design
search, the prescribed candidate grid used states $-q,0,q$ and scales
$1,2,3,4$. Of its $\binom{12}{8}=495$ eight-row subsets, 486 were exactly
nonsingular. The selected design was chosen deterministically by the
prescribed preferences and a numerical conditioning tie-break; nonsingularity
was certified exactly afterward. Each one-row deletion of $B$ has rank
seven. None of these search counts is needed to prove the inverse once
the single exact matrix is supplied.

\section{Quantitative local stability at the scalar witness}
\label{sec:stability}
For vectors use the maximum norm and for matrices its induced row-sum
norm. Define $K=\norm{B^{-1}}_\infty$. Every component of $\Phi_*$ is a
polynomial of ordinary coordinate degree at most four. Hence write
\begin{equation}\label{eq:derivative_poly}
 (D\Phi_*(z)-B)_{ja}
 =\sum_{1\leq|\alpha|\leq3} c_{ja,\alpha}z^\alpha,
\end{equation}
where $\alpha=(\alpha_1,\ldots,\alpha_8)$ is a nonnegative integer
multiindex, $|\alpha|=\sum_a\alpha_a$, and
$z^\alpha=\prod_a z_a^{\alpha_a}$. Put
\begin{equation}\label{eq:Delta}
 a_{j,p}=\sum_{a=1}^8\sum_{|\alpha|=p}|c_{ja,\alpha}|,
 \qquad
 \Delta(R)=\max_{1\leq j\leq8}\sum_{p=1}^3a_{j,p}R^p\quad(R>0).
\end{equation}
This is an upper majorant, not a claim to equal the supremum of the
Jacobian difference. Coefficientwise triangle inequalities give
\begin{equation}\label{eq:majorant}
 \sup_{\norm{z}_\infty\leq R}\norm{D\Phi_*(z)-B}_\infty\leq\Delta(R).
\end{equation}
The exact row sums are recorded in \cref{tab:majorant}; the full
polynomial coefficient specification appears in \cref{app:polynomial}.

\begin{theorem}[Exact local inverse-Lipschitz certificate]
\label{thm:stability}
At the fixed witness and design \eqref{eq:design}, set
\begin{align}\label{eq:stabilityvalues}
 K&=\frac{23812}{2151},& R_*&=\frac1{2048},\nonumber\\
 \Delta_*&=\frac{2643019857}{68719476736},&
 \eta_*:=K\Delta_*&=\frac{1748210800969}{4105988734976}<\frac12,\nonumber\\
 c_*&=\frac{21220001406063}{409087045009408},&
 C_*&=\frac{409087045009408}{21220001406063}.
\end{align}
For all $z,z'$ in the closed cube
$\mathcal C_*=[-R_*,R_*]^8$,
\begin{equation}\label{eq:stability}
 \norm{\Phi_*(z)-\Phi_*(z')}_\infty\geq
 c_*\norm{z-z'}_\infty,
 \qquad
 \norm{z-z'}_\infty\leq C_*\norm{\Phi_*(z)-\Phi_*(z')}_\infty.
\end{equation}
In particular $\Phi_*$ is injective on that cube, and the inverse on its
image is Lipschitz. A simpler valid upper bound for the inverse constant
is $2K=47624/2151$.
\end{theorem}
\begin{proof}
The exact inverse rows in \cref{app:inverse} give the displayed $K$.
The polynomial coefficients in \cref{app:polynomial}, differentiated as
in \eqref{eq:derivative_poly}, give \cref{tab:majorant}. Evaluating
\eqref{eq:Delta} at $R_*$ yields $\Delta_*$. Exact multiplication gives
$\eta_*$ and
\[
 \frac12-\eta_*=\frac{304783566519}{4105988734976}>0.
\]
Let $v=z-z'$ and $\gamma(t)=z'+tv$ for $0\leq t\leq1$.
Convexity puts the entire segment in $\mathcal C_*$. The fundamental
theorem of calculus, preconditioned by the fixed matrix $B^{-1}$, gives
\[
 B^{-1}(\Phi_*(z)-\Phi_*(z'))
 =v+\int_0^1B^{-1}(D\Phi_*(\gamma(t))-B)v\,dt.
\]
The integral term has norm at most $K\Delta_*\norm{v}_\infty$ by
\eqref{eq:majorant}. The reverse triangle inequality around the fixed
vector $v$, followed by submultiplicativity, gives
\[
 K\norm{\Phi_*(z)-\Phi_*(z')}_\infty
 \geq(1-K\Delta_*)\norm{v}_\infty.
\]
Thus the lower constant is $1/K-\Delta_*=c_*>0$ and its reciprocal is
$C_*$. The inequality $\eta_*<1/2$ gives $C_*\leq2K$. This argument
does not take a lower bound inside the norm of a vector integral.
\end{proof}

The radius is the largest passing value among the sixteen prescribed
candidates $2^{-1},\ldots,2^{-16}$ for the test $K\Delta(R)\leq1/2$;
it is not asserted to be the largest possible invertibility radius.
\Cref{tab:radii} gives the complete rational search table. If $z$ and a
decoded $\widehat z$ both belong to $\mathcal C_*$ and
$\Phi_*(\widehat z)=\Phi_*(z)+e$, then
$\norm{\widehat z-z}_\infty\leq C_*\norm e_\infty$.
Existence of such a preimage for arbitrary noisy data is an additional
condition. The certificate supplies no parameter-uniform neighborhood.

\begin{proposition}[Common-scale determinant and amplification laws]
\label{prop:scales}
Let $h\ne0$ multiply every scale in \eqref{eq:design}, and let $B(h)$
be the resulting zero-coordinate Jacobian. With
$\mathbf d=(1,1,2,3,3,4,4,4)$,
\begin{equation}\label{eq:scaleB}
 B(h)=B\diag(h^{d_1},\ldots,h^{d_8}),\qquad
 \det B(h)=h^{22}\det B.
\end{equation}
Its exact inverse amplification in the maximum norm is
\begin{equation}\label{eq:scalenorm}
 \norm{B(h)^{-1}}_\infty=\max\left\{
 \frac{4547}{717}|h|^{-1},\frac{23812}{2151}|h|^{-2},
 \frac{824}{99}|h|^{-3},\frac{160}{33}|h|^{-4}\right\}.
\end{equation}
For $0<|h|\leq1$ this lies between
$(160/33)|h|^{-4}$ and $(23812/2151)|h|^{-4}$.
\end{proposition}
\begin{proof}
Equation \eqref{eq:zero_derivative} multiplies column $a$ by $h^{d_a}$.
The degrees sum to $1+1+2+3+3+4+4+4=22$, proving the determinant law.
Inverting \eqref{eq:scaleB} multiplies row $a$ of $B^{-1}$ by
$h^{-d_a}$. Take the exact row sums in \cref{eq:invrows} and maximize
within each degree to obtain \eqref{eq:scalenorm}. The last bounds follow
because $|h|^{-p}\leq|h|^{-4}$ for $p\leq4$ and the largest unscaled
row sum is $K$.
\end{proof}
This is a law for the inverse linearization, not for the product of
forward and inverse norms called a condition number. The nonlinear
variation majorant changes with the design, so the same cube and constant
in \cref{thm:stability} are not asserted after rescaling.

\subsection{Approximate recovery from actual responses}
\label{sec:actual}
The following result uses the fixed scalar witness and all eight designed
experiments. Write $\Phi=\Phi_*$, $R=R_*$ and $Q_R=\mathcal C_*$, with
$B,K,\Delta_*$ from \cref{thm:stability}. All matrix infinity norms below
are induced row-sum norms. In particular, put
\[
 \kappa=K\Delta_*=\eta_*<\tfrac12,\qquad
 C=\frac{K}{1-\kappa}=C_*.
\]

\begin{theorem}[Approximate Hall recovery from eight actual responses]
\label{thm:actual}
Use $V_1(y)=\tanh y$, $V_2(y)=1/2$ and exactly the pairs
$(y_j,r_j)$ in \eqref{eq:design}. Let $X:[0,T]\to\R^2$ be continuous,
of bounded variation, with $X_0=0$ and
$\operatorname{TV}(X^1)+\operatorname{TV}(X^2)\leq M$, where $M\geq0$.
Let $z(X)$ be its Hall coordinates through degree four in
\eqref{eq:Hall8}. Set
\[
 D_h=\diag(h,h,h^2,h^3,h^3,h^4,h^4,h^4),\qquad
 (c_1,c_2,c_3,c_4)=(1,1,1,13/12),\qquad A=3402.
\]
Choose $0<h\leq1$ such that $c_m(hM)^m\leq R$ for $m=1,\ldots,4$.
A sufficient choice for $M>0$ is $h\leq\min(1,R/M)$; if $M=0$, any
$0<h\leq1$ is allowed. The full CDE solutions exist, and their centered
observations
\[
 b_h(X)_j=Y_T^{y_j,hr_j}(X)-y_j
\]
satisfy $D_hz(X)\in Q_R$ and
\begin{equation}\label{eq:actualremainder}
 \norm{b_h(X)-\Phi(D_hz(X))}_\infty\leq Ah^5M^5.
\end{equation}
For any $b\in\R^8$, let $\Pi_{Q_R}$ be coordinatewise clipping onto $Q_R$
and define
\begin{equation}\label{eq:clippeddecoder}
 T_b(u)=\Pi_{Q_R}\bigl[u+B^{-1}(b-\Phi(u))\bigr],\qquad u\in Q_R.
\end{equation}
This is a contraction with constant at most $\kappa$. Its unique fixed
point $G(b)$ exists for every $b$, satisfies $G(\Phi(u))=u$ for $u\in Q_R$,
and obeys
\[
 \norm{G(b)-G(b')}_\infty\leq C\norm{b-b'}_\infty.
\]
If $b=b_h(X)+e$, with $\norm e_\infty\leq\delta$, then
$\widehat z_h=D_h^{-1}G(b)$ satisfies
\begin{equation}\label{eq:actualerror}
 \norm{\widehat z_h-z(X)}_\infty
 \leq C\bigl(AM^5h+\delta h^{-4}\bigr).
\end{equation}
More generally, for $u_n\in Q_R$, put
$r_n=\norm{u_n-T_b(u_n)}_\infty$. Then
\begin{equation}\label{eq:iterationerror}
 \norm{D_h^{-1}u_n-z(X)}_\infty
 \leq C\bigl(AM^5h+\delta h^{-4}\bigr)
       +\frac{h^{-4}r_n}{1-\kappa}.
\end{equation}
Here $\delta$ bounds response-observation error and is unrelated to the
financial observation-noise parameter.
\end{theorem}
\begin{proof}
The truncation in \eqref{eq:truncated} is by tensor word length.
In \eqref{eq:polyfull}, a Hall monomial $z^\alpha$ has weight
$\mathrm{wt}(\alpha)=\sum_a d_a\alpha_a$, with
$\mathbf d=(1,1,2,3,3,4,4,4)$, and carries the scale factor
$r^{\mathrm{wt}(\alpha)}$. All weights are at most four. Thus replacing
every design scale by $hr_j$ gives the full polynomial identity
\begin{equation}\label{eq:fullpolyscaling}
 \Phi_h(z)=\Phi(D_hz),
\end{equation}
including every nonlinear term.

Let $\mu=|dX^1|+|dX^2|$ and let $|\,\cdot\,|_1$ denote the sum of
absolute tensor-word coefficients. The signature simplex formula gives
$|S_m(X)|_1\leq M^m/m!$, and concatenation is submultiplicative in this
norm. The finite tensor logarithm therefore gives
\[
 |\ell_m(X)|_1\leq
 M^m\sum_{k=1}^m\frac1k
 \sum_{\substack{j_1+\cdots+j_k=m\\j_a\geq1}}
 \frac1{j_1!\cdots j_k!}=c_mM^m\qquad(m\leq4).
\]
In the Hall order \eqref{eq:Hall8}, the word-expansion submatrices on
rows $(1,2)$, $(12)$, $(112,122)$ and $(1112,1122,1222)$ are respectively
identity matrices. Each Hall coordinate is consequently the corresponding
word coefficient of $\ell_m$. Hence
$\norm{z_m(X)}_\infty\leq c_mM^m$, proving the required scaled domain
condition without a restriction $z(X)\in Q_R$. Since $R=1/2048$,
$c_mR^m\leq R$ for all four degrees, which proves the simpler sufficient
scale condition.

Both fields are globally bounded and globally Lipschitz. At any finite
scale $a$, the CDE integral map is a contraction on successive intervals
of sufficiently small $\mu$-mass. Such a finite partition exists because
$\mu$ is finite and nonatomic. Moreover
$\sup_t|Y_t-y|\leq |a|M$, so these local solutions extend uniquely through
$T$. Set $s=\tanh y$, $v_1(s)=s$, $v_2(s)=1/2$. The word convention
\eqref{eq:Gamma} gives
\[
 P_i=v_i,\qquad P_{iv}=v_i(s)(1-s^2)P_v'(s),\qquad
 (\Gamma(w)\Id)(y)=P_w(\tanh y).
\]
For the polynomial coefficient norm $a(P)=\sum_k|[s^k]P|$,
$\deg P_v\leq2n-1$ when $|v|=n$, and
$a(P_{iv})\leq2(2n-1)a(P_v)$. Starting from
$\max_i a(P_i)=1$, four steps give the exact upper bound
\[
 \max_{|w|=5}\sup_y|(\Gamma(w)\Id)(y)|
 \leq 2^4(1\cdot3\cdot5\cdot7)=:K_5=1680.
\]
Apply the finite BV chain-rule expansion from \cref{prop:jets} through
word length four. At scale $a$, its remainder is
\[
 a^5\sum_{i_1,\ldots,i_5}
 \int_{0<t_1<\cdots<t_5<T}
 P_{i_1\cdots i_5}(\tanh Y_{t_1})
 \,dX^{i_1}_{t_1}\cdots dX^{i_5}_{t_5}.
\]
Taking absolute values and summing over words gives the product measure
$\mu^{\otimes5}$, whose ordered simplex has mass $\mu([0,T])^5/5!$.
There is no additional word-count factor. Thus the remainder is bounded
by $K_5|a|^5M^5/5!$. Since $\max_j|r_j|=3$,
$A=3^5K_5/5!=3402$, proving \eqref{eq:actualremainder}.
This finite expansion needs no infinite-series convergence. Its bound is
uniform in the initial state. The constants $A,R,K,\kappa,C$ depend only on
the fixed witness, design and certificate, and not on $h$, $M$, $X$ or $T$.

By \eqref{eq:majorant}, the derivative of
$u\mapsto u+B^{-1}(b-\Phi(u))$ has induced norm at most
$K\Delta_*=\kappa$ on the convex cube. Clipping is nonexpansive in the
maximum norm, so $T_b$ is a contraction of the complete space $Q_R$.
The contraction theorem supplies its unique fixed point for every $b$.
If $b=\Phi(u)$, then $T_b(u)=u$, even on the boundary, proving exactness
there. Nonexpansiveness also gives
\[
 \norm{T_b(u)-T_{b'}(v)}_\infty
 \leq\kappa\norm{u-v}_\infty+K\norm{b-b'}_\infty.
\]
Apply this to the two fixed points and rearrange to obtain the claimed
Lipschitz bound. With $u=D_hz(X)\in Q_R$, it follows that
\[
 \norm{G(b)-u}_\infty\leq C(Ah^5M^5+\delta).
\]
Finally $\norm{D_h^{-1}}_\infty=h^{-4}$ for $0<h\leq1$, proving
\eqref{eq:actualerror}. The contraction estimate also implies
$\norm{u_n-G(b)}_\infty\leq r_n/(1-\kappa)$, which gives
\eqref{eq:iterationerror}.
\end{proof}

The map $G$ extends the inverse on $\Phi(Q_R)$ to all observation vectors;
it is not a global inverse of $\Phi$ or a least-squares decoder. A clipped
fixed point may have zero fixed-point residual and nonzero response-fit
residual $\norm{\Phi(u)-b}_\infty$. To keep the iteration contribution
below an unscaled coordinate tolerance $\tau$, the mathematical stopping
criterion is $r_n\leq(1-\kappa)h^4\tau$. A numerically computed residual
also contains roundoff and is not by itself a certified exact residual.
The result gives approximate recovery from the eight designed actual
responses, not exact recovery at fixed nonzero scale or recovery from one
zero-initialized scalar reservoir.

For $M=1$ and $h=1/2048$, the noiseless term in
\eqref{eq:actualerror} is approximately $32.0239$, and the coefficient
of absolute response error $\delta$ is approximately
$3.39149\times10^{14}$. These are conservative upper bounds, not observed
errors or evidence of numerical failure. The asymptotic consistency
statement does not by itself certify useful finite-precision accuracy at
a specified scale; response error, iteration error and roundoff still
require control.

For deterministic axis-aligned piecewise-linear paths, the actual responses
can be evaluated using the exact segment flows
\[
 F_a(y)=\operatorname{arsinh}(e^a\sinh y),\qquad H_b(y)=y+b/2.
\]
They satisfy $\partial_aF_a(y)=\tanh(F_a(y))$ and
$\partial_bH_b(y)=1/2$, with identity initial maps. Compose them in path
order, using signed increments multiplied by $hr_j$; the last segment acts
outermost on the state. The reference signature instead concatenates the
segment tensor exponentials on the right, as in \eqref{eq:signature}.
The validation command in \cref{app:reproduce} compares these independently
constructed responses and Hall coordinates at two arbitrary precisions.

\section{Lifting the scalar witness in state dimension}
\label{sec:lift}
\begin{corollary}[Inactive-coordinate scalar embedding]\label{cor:lift}
For every $D\geq1$, the two-channel shallow componentwise-tanh family on
$\R^D$ has a witness with homogeneous free-Lie injectivity through degree
four obtained from \eqref{eq:witness}. At that witness the mixed-degree
kernel relation \eqref{eq:kernel8} also persists. First-coordinate
observations at embedded initial states reproduce \eqref{eq:Phi} exactly,
including the stability certificate. Fixed lifted homogeneous minors give
an open dense full-measure certificate set for separate homogeneous
injectivity in the full parameter space $\R^{2(D^2+D)}$.
\end{corollary}
\begin{proof}
Define $E(f)(x)=f(x_1)e_1$, where $e_1$ is the first coordinate vector.
Differentiating shows $[E(f),E(g)]=E(fg'-f'g)$. Thus $E$ is an injective
Lie homomorphism. Choose $A_1$ with sole nonzero entry $(A_1)_{11}=1$,
$b_1=0$, $A_2=0$, and $b_2=qe_1$. The fields are exactly
$E(\tanh)$ and $E(1/2)$, because $\tanh(0)=0$ in inactive coordinates.
For a nonempty word, induction using
$F_{iv}=DF_v\,V_i$ shows that its response equals its scalar response
in coordinate one and zero in the other coordinates, globally in $x$.
Consequently every term of the truncated response has this property.
Initial states $y e_1$ therefore reproduce the scalar centered map, not
just its derivative. The homogeneous certificates, mixed-degree relation,
and quantitative constants at this witness follow.

For arbitrary full weights and biases, evaluate the same homogeneous
Hall fields at the same embedded states and select coordinate one.
Their determinant product is analytic in every parameter and equals
\eqref{eq:dets} at the embedding. \Cref{lem:analytic} proves generic
separate homogeneous injectivity. No stability constant is propagated
to that generic parameter set.
\end{proof}
The embedding is a corollary of an existing witness. It is not a statement
that increasing state dimension improves every reservoir or every
conditioning measure.

\section{A general sufficient-width shallow-tanh theorem}
\label{sec:width}
For the full architecture \eqref{eq:architecture}, we now seek
independence of all nonempty word-response fields through length $L$.
This is a different linear-family requirement from the Hall evaluation
rank in \cref{thm:hallrank}. No tree-indexed independence condition is
introduced.

\subsection{A transfer lemma for finite word-response minors}
\begin{lemma}[Normalized linear-to-tanh transfer]\label{lem:transfer}
Let $W_i(x)=A_ix$ be linear fields on $\R^D$. Suppose a finite family
of nonempty words $w_1,\ldots,w_s$ has a nonsingular square matrix under
fixed scalar point-coordinate evaluations $\lambda_1,\ldots,\lambda_s$:
\[
 M_0=(\lambda_j(\Gamma_W(w_a)\Id))_{j,a}.
\]
For $V_i^\epsilon(x)=\tanh(\epsilon A_ix)$, the same evaluated word
matrix is nonsingular for all sufficiently small nonzero $\epsilon$.
The conclusion asserts existence of a punctured neighborhood, not a
certified explicit value of $\epsilon$.
\end{lemma}
\begin{proof}
Introduce the auxiliary fields
\begin{equation}\label{eq:normalized}
 \widetilde V_i^\epsilon(x)=
 \begin{cases}\epsilon^{-1}\tanh(\epsilon A_ix),&\epsilon\ne0,\\
 A_ix,&\epsilon=0.
 \end{cases}
\end{equation}
The Taylor series of $\tanh s$ at zero is odd, begins with $s$, and has
positive convergence radius. After division by $\epsilon$ it is a
convergent series in even powers of $\epsilon$, with polynomial
coefficients in $x$. On a neighborhood of each of the finitely many
evaluation states it defines a jointly real-analytic extension at
$\epsilon=0$. All finite state derivatives required by the chosen words
are therefore real analytic there. Induction using
$F_{iv}=DF_v\,V_i$ proves that
\[
 A(\epsilon)_{j,a}:=\lambda_j
 (\Gamma_{\widetilde V^\epsilon}(w_a)\Id)
\]
is analytic near zero and $A(0)=M_0$.

For $\epsilon\ne0$, the scalar factor is independent of state, so
$\Dop_{V_i^\epsilon}=\epsilon\Dop_{\widetilde V_i^\epsilon}$.
Every word therefore obeys the exact factorization
\begin{equation}\label{eq:epsfactor}
 \Gamma_{V^\epsilon}(w)\Id
 =\epsilon^{|w|}\Gamma_{\widetilde V^\epsilon}(w)\Id.
\end{equation}
The raw evaluated determinant is
$\epsilon^{\sum_a|w_a|}\det A(\epsilon)$.
Since $\det A(0)=\det M_0\ne0$, continuity makes
$\det A(\epsilon)\ne0$ for $|\epsilon|$ sufficiently small. Removing
zero also makes the prefactor nonzero, proving the claim.
\end{proof}

The raw fields at $\epsilon=0$ are zero; their evaluated word matrix is
not $M_0$. Only the normalized auxiliary matrix has that limit. The
argument uses a local determinant limit, without an all-$\epsilon$
triangularity or determinant formula.

\subsection{The classical tensor system in chronological convention}
Let
\begin{equation}\label{eq:D0}
 D_0=D_0(d,L)=\sum_{m=0}^Ld^m,\qquad
 Q_{d,L}=\sum_{m=1}^Ld^m=D_0-1.
\end{equation}
Use the truncated tensor-word space as $\R^{D_0}$, with basis $e_w$
ordered first by word length and then lexicographically, including
$e_{\varnothing}$. Define right-concatenation matrices
\begin{equation}\label{eq:Rconcat}
 R_i e_w=\begin{cases}e_{wi},&|w|<L,\\0,&|w|=L,
 \end{cases}
 \qquad W_i(x)=R_ix,\qquad x_0=e_{\varnothing}.
\end{equation}
This is the classical truncated tensor system; related formulations occur
in \cite{FrizVictoir,AkyildirimAnomaly}. Its action orientation must be
checked against \eqref{eq:Gamma}.

For a linear map $A$, $\Dop_{W_i}(Ax)=AR_ix$. Induction therefore gives
\begin{equation}\label{eq:linearwords}
 \Gamma_W(i_1\cdots i_m)\Id(x)=R_{i_m}\cdots R_{i_1}x,
 \qquad
 \Gamma_W(w)\Id(x_0)=e_w\quad(1\leq|w|\leq L).
\end{equation}
The rightmost matrix acts first and appends the earliest letter first.
Thus no word reversal is inserted into the chronological signature.
Select all nonempty-word coordinates as rows and all nonempty words
through $L$ as columns. Equation \eqref{eq:linearwords} gives the exact
square matrix
\begin{equation}\label{eq:A0}
 A_0=(\Gamma_W(w)\Id(x_0)_u)_{u,w}=I_{Q_{d,L}}.
\end{equation}

Products of $L+1$ right-concatenation matrices vanish. For a segment with
increment $a$, the linear flow at $x_0$ is
$\exp(\sum_i a_iR_i)x_0=\sum_{m=0}^L a^{\otimes m}/m!$.
Two segment flows apply the second exponential on the left of the first
matrix exponential; on tensor states this appends the second tensor
signature on the right of the first, producing the Chen product in
\eqref{eq:signature}. These identities also explain the exact checks in
\cref{app:reproduce}.

\subsection{The strongest conclusion and full parameter genericity}
\begin{theorem}[Finite-degree word-response independence at sufficient width]
\label{thm:width}
For every finite $d,L\geq1$ and every integer $D\geq D_0(d,L)$, there
exists an open dense full-Lebesgue-measure subset
$\mathcal G_{d,L,D}^{\rm cert}\subset\mathcal P_{d,D}$ such that the
finite family of vector-valued functions on $\R^D$
\begin{equation}\label{eq:wordfamily}
 \{\Gamma_\theta(w)\Id:1\leq|w|\leq L\}
\end{equation}
is jointly linearly independent for every
$\theta\in\mathcal G_{d,L,D}^{\rm cert}$. Consequently
$\rho_\theta:\Lie_d^{\leq L}\to\mathcal X(\R^D)$ is injective as a
linear map on the truncated free-Lie vector space and on every homogeneous
component separately. The certificate uses a fixed initial state and
$Q_{d,L}$ fixed output-coordinate evaluations of the word-response fields.
\end{theorem}
\begin{proof}
First take $D=D_0$. Apply \cref{lem:transfer} to \eqref{eq:Rconcat} and
the identity matrix \eqref{eq:A0}. Equivalently, let
$V_i^\epsilon=\tanh(\epsilon R_ix)$ and define the normalized matrix
at $x_0$ with the nonempty-word rows and columns. It is real analytic
near zero and has value $I_{Q_{d,L}}$ there. In this case the raw
determinant is exactly
\begin{equation}\label{eq:rawdet}
 \epsilon^{\kappa_{d,L}}\det A(\epsilon),\qquad
 \kappa_{d,L}=\sum_{m=1}^L m d^m,
 \qquad \det A(0)=1.
\end{equation}
Choose existentially one sufficiently small nonzero $\epsilon$.
The corresponding weights $A_i=\epsilon R_i$, biases $b_i=0$, give
a nonzero evaluated word determinant in the actual architecture.
If a linear combination of the word-response fields vanished globally,
evaluating the chosen coordinates at $x_0$ would put its coefficients
in the kernel of this nonsingular matrix. Thus they are independent.

For an arbitrary parameter $\theta\in\mathcal P_{d,D_0}$ define the
fixed determinant
\begin{equation}\label{eq:genericminor}
 \Delta_{d,L,D_0}(\theta)=
 \det\bigl(\Gamma_\theta(w)\Id(x_0)_u\bigr)_{u,w}.
\end{equation}
The map $(\theta,x)\mapsto\tanh(A_ix+b_i)$ is jointly real analytic
for all real parameters and states. Recursing by
$F_{iv}^\theta=D_xF_v^\theta V_i^\theta$ proves joint analyticity
of each finite word-response field. Evaluation and a finite determinant
preserve analyticity. Thus \eqref{eq:genericminor} is analytic on the
entire connected Euclidean parameter space and is not identically zero
by the constructed point. Its nonzero set has the claimed topological
and measure properties by \cref{lem:analytic}.

For $D>D_0$, embed $x_0$ in the first $D_0$ coordinates, put the
constructed matrices in the upper-left block, and set all other weights
and biases to zero. Writing $x=(x^{\rm act},x^{\rm inact})$, the fields
are $(V_i^{\rm act}(x^{\rm act}),0)$ globally, independent of the
inactive variables. The recursion for word responses preserves this
form, so the selected active-coordinate minor is exactly the smaller
one. Define its determinant for arbitrary full $D$-dimensional weights
and biases using the same embedded state and rows. The preceding
analyticity and nonzero-witness argument applies to every entry of
the full parameter vector, proving the result for this $D$.

Finally let $0\ne H\in\Lie_d^{\leq L}$. Its tensor expansion is a
nonzero real combination of the word basis through $L$. By
\eqref{eq:Lieoperator}, its represented field is the same combination
of \eqref{eq:wordfamily}. Word-response independence prevents it from
vanishing. This proves joint Lie injectivity, and restriction proves
the homogeneous conclusions.
\end{proof}

The nonzero locus is a certificate set; a parameter where this particular
minor vanishes might still admit another independence certificate.
An absolutely continuous parameter law assigns probability one to the
certificate set. This does not cover a randomization supported on an
exceptional lower-dimensional set.

The theorem concerns lengths at most $L$ and constant real linear
relations among word-response fields. It does not prove tree-level
independence, all-degree independence, or factorization of the tanh
action through the free step-$L$ nilpotent quotient. The explicit
dimension $D_0$ is a conservative sufficient width, with no lower-bound,
minimality, compressed-width, or monotonic-performance claim. The proof
provides no certified explicit nonzero $\epsilon$, conditioning estimate,
parameter-uniform stability constant, or quantitative inverse theorem
for these large witnesses. A field-evaluation determinant at one state
is not reconstruction of a path from one terminal state. Full-CDE
reconstruction, discrete-time transfer, and finance consequences remain
outside the theorem.

\begin{corollary}[Fixed zero initialization]\label{cor:zero}
Let $d,L\geq1$ and $D\geq D_0(d,L)=\sum_{m=0}^L d^m$. In the full
biased architecture \eqref{eq:architecture}, fix the initial state $h_0=0$.
For Lebesgue-almost-every $\theta\in\mathcal P_{d,D}$, the vectors
\[
 \{(\Gamma_\theta(w)\Id)(0):1\leq|w|\leq L\}
\]
are jointly linearly independent. Consequently the finite-degree
word-response fields are jointly independent, and $\rho_\theta$ is
injective on $\Lie_d^{\leq L}$ and on each homogeneous component, as in
\cref{thm:width}.
\end{corollary}
\begin{proof}
Put $c=e_{\varnothing}$ in the first $D_0$ coordinates of $\R^D$,
with zero inactive coordinates. For any fields
$V_i(x)=\tanh(A_ix+b_i)$, let $\tau_c(z)=z+c$ and set
$W_i(z)=\tanh(A_iz+b_i+A_ic)$.
For every coordinate $k$,
$(W_i(z))_k=\tanh((A_i(z+c)+b_i)_k)=(V_i(z+c))_k$.
Thus translation stays in the same architecture, with weights $A_i$
and biases $b_i+A_ic$.

The chain rule gives
$\Dop_{W_i}(f\circ\tau_c)=(\Dop_{V_i}f)\circ\tau_c$.
Prepending letters and using the outermost-first convention
\eqref{eq:Gamma} proves by induction that
\[
 \Gamma_W(w)(f\circ\tau_c)=(\Gamma_V(w)f)\circ\tau_c.
\]
Since $\Id\circ\tau_c=\Id+c$, take $f=\Id$. Every nonempty word
operator kills the constant vector $c$: its rightmost first-order
operator differentiates that constant to zero. Hence
\[
 (\Gamma_W(w)\Id)(0)=(\Gamma_V(w)\Id)(c),
 \qquad 1\leq|w|\leq L.
\]

Use precisely the $Q_{d,L}\times Q_{d,L}$ minor in \eqref{eq:A0}:
rows are the active coordinates indexed by nonempty words $u$ through
$L$, and columns are the nonempty words $w$ through $L$, both in
degree-then-lexicographic order. No empty-word or inactive row is used.
The preceding identity preserves this minor entry by entry.
At $D=D_0$, the proof of \cref{thm:width} gives its raw determinant
$\epsilon^{\kappa_{d,L}}\det A(\epsilon)$, where $A(0)=I_{Q_{d,L}}$.
Continuity therefore permits a finite, sufficiently small nonzero
$\epsilon$ with nonzero determinant. Translating these actual tanh
fields gives admissible weights $A_i=\epsilon R_i$ and biases
$b_i=\epsilon R_ic$. The normalized limit at $\epsilon=0$ is only the
existence argument, not the witness. For $D>D_0$, embed these weights
and biases in the active block and set all other entries to zero;
the same minor is unchanged.

With the rows and columns just fixed, define on the full parameter space
\[
 \Delta^0_{d,L,D}(\theta)
 =\det\bigl((\Gamma_\theta(w)\Id)(0)_u\bigr)_{u,w}.
\]
Finite differentiation of the jointly real-analytic tanh fields makes
this a real-analytic function of all entries of $(A_i,b_i)_{i=1}^d$,
with the state held at zero. The translated witness makes it nontrivial,
so \cref{lem:analytic} gives a null zero set. Nonsingularity proves the
stated independence; the tensor embedding and \eqref{eq:Lieoperator}
give the Lie conclusions.

The empty word is excluded throughout: $\Gamma(1)\Id(0)=0$, so adding
its column would destroy injectivity of the evaluation map. Its signature
coefficient is the known value $S_0(X)=1$ and is treated separately.
\end{proof}

\section{Relationship between the architecture results}
\label{sec:comparison}
The fixed-state rank theorem, the scalar certificates, and the word
theorem address different observations and different families.
For $d=2,L=4$, \cref{thm:hallrank,cor:hallinverse} give full Hall rank
and local recovery of eight truncated Hall coordinates from one
zero-initialized vector at width $\nu_{2,4}=8$. The scalar witness at
$D=1$ instead has joint function rank seven, despite homogeneous
injectivity through degree four. Its eight selected state-scale
observations admit the explicit inverse and stability bound of
\cref{sec:eight,sec:stability}, and its actual CDE responses admit
the approximate decoder of \cref{thm:actual}.

At the larger width $D_0(2,4)=31$, \cref{thm:width} proves independence
of all $Q_{2,4}=30$ nonempty word-response fields.
Corollary~\ref{cor:zero} also gives their independent evaluations at zero.
The Hall-width construction does not assert independence of all these
word directions in eight linear features.

The reconstruction implications are
\begin{align*}
 \text{full fixed-state Hall rank}
 &\ \Longrightarrow\ \text{one-vector truncated local inverse},\\
 \text{homogeneous field injectivity}
 &\ \Longrightarrow\ \text{ideal graded Hall recovery},\\
 \text{homogeneous field injectivity}
 &\ \Longrightarrow\ \text{finite state-scale truncated inverse}.
\end{align*}
Independent Hall evaluations imply joint field independence and hence
homogeneous injectivity. Function-level independence alone need not give
independent evaluations at one state. Word-response field independence
implies joint Hall field independence by the tensor embedding, but is a
larger linear-family requirement.

The scalar inverse becomes quantitative only through its separate exact
Jacobian and variation bounds. Its extension to actual responses also
uses a uniform remainder and a clipped contraction on the certified
cube. None of those constants, scale exponents, or remainder estimates
is transferred to the Hall-width or word-width construction.

\section{Financial motivation and the information constraint}
\label{sec:finance}
The finance statements in this section are classical probability and
projection facts. They explain why predictable history features are useful;
they make no claim that a finite-degree identifiability theorem determines a
trained model's financial risk.

\subsection{Predictable martingale controls for a discrete pricing target}
Let $(\Omega,\mathcal F,\mathbb Q)$ be a probability space, and let
$0=t_0<\cdots<t_N=T$ be a simulation grid. Let $G_N\in L^2(\mathbb Q)$ be
the discounted payoff computed on that grid. The target is
$p_N=\E_{\mathbb Q}G_N$. Write $\mathcal F_k^{\rm sim}$ for simulator
information through time $t_k$. A simulator martingale innovation
$\Delta M_{k+1}\in\R^b$ satisfies
$\E[\Delta M_{k+1}\mid\mathcal F_k^{\rm sim}]=0$, where $b$ is the number
of innovation components. Brownian simulation increments are an example.
They need not be increments of traded assets.

Let $\mathcal T$ contain the reservoir draw, preprocessing, training data,
fitting, and validation selection. Evaluation paths are independent of
$\mathcal T$. A frozen predictor $Z_k\in\R^b$ is measurable with respect
to $\mathcal T\vee\mathcal F_k^{\rm sim}$ and uses only a path prefix. Its
control is
\begin{equation}\label{eq:control}
 C(Z)=\sum_{k=0}^{N-1}Z_k\cdot\Delta M_{k+1}.
\end{equation}
The innovation assumption below is stated conditional on the enlarged
filtration so that the role of sample splitting is explicit.

\begin{proposition}[Sample-split unbiasedness]\label{prop:unbiased}
Let $G_N\in L^2(\mathbb Q)$ and $p_N=\E_{\mathbb Q}G_N$. Assume that
each evaluation path has the prescribed simulation law and is independent
of the training sigma-field $\mathcal T$. For every $k=0,\ldots,N-1$,
assume that $Z_k$ is $\mathcal T\vee\mathcal F_k^{\rm sim}$-measurable,
that $\Delta M_{k+1}$ is integrable with
$\E[\Delta M_{k+1}\mid\mathcal T\vee\mathcal F_k^{\rm sim}]=0$, and
that every product in \eqref{eq:control} is integrable. For $P$ evaluation
paths that are independent conditional on $\mathcal T$, the estimator
\[
 \widehat p_N=\frac1P\sum_{j=1}^P\bigl(G_N^{(j)}-C^{(j)}(Z)\bigr)
\]
satisfies $\E[\widehat p_N\mid\mathcal T]=p_N$.
\end{proposition}
\begin{proof}
Predictability and conditional centering give
\[
 \E[Z_k\cdot\Delta M_{k+1}\mid\mathcal T]
 =\E\!\left[Z_k\cdot
 \E[\Delta M_{k+1}\mid\mathcal T\vee\mathcal F_k^{\rm sim}]
 \mid\mathcal T\right]=0.
\]
Sum over $k$. Independence of the evaluation path from $\mathcal T$ gives
$\E[G_N\mid\mathcal T]=p_N$, and linearity proves the assertion.
\end{proof}

This proves unbiasedness for the discretized target only. It also explains
two implementation details: a fitted intercept is not subtracted as a
martingale control, and training standardization must be converted back to
coefficients of the original predictable control features before evaluation.
Centering controls by an evaluation-dependent fit is not part of this
argument. Martingale control variates in stochastic-volatility pricing have
well-established precedents, including Fouque and Han \cite{FouqueHan}.

\begin{proposition}[Finite-dimensional control projection]\label{prop:projection}
Let $G_N\in L^2(\mathbb Q)$, put $p_N=\E_{\mathbb Q}G_N$, and let
$Q\in L^2(\mathbb Q;\R^p)$ be a centered vector of terminal controls.
Set $g=G_N-p_N$, $\Sigma=\E[QQ^\top]$, and $c=\E[Qg]$.
Let $\Sigma^\dagger$ be the Moore--Penrose inverse. Then
$c\in\range(\Sigma)$, and for every $\beta\in\R^p$,
\begin{align}\label{eq:projection}
 \Var(G_N-\beta^\top Q)
 &=\Var(G_N)-c^\top\Sigma^\dagger c\nonumber\\
 &\quad +(\beta-\Sigma^\dagger c)^\top
 \Sigma(\beta-\Sigma^\dagger c).
\end{align}
The minimizing coefficients are $\Sigma^\dagger c+\ker\Sigma$.
\end{proposition}
\begin{proof}
For $v\in\ker\Sigma$, $\E[(v^\top Q)^2]=0$, so $v^\top Q=0$ almost
surely and $v^\top c=0$. Symmetry of $\Sigma$ gives
$c\in(\ker\Sigma)^\perp=\range(\Sigma)$. Expanding the variance gives
$\E[g^2]-2\beta^\top c+\beta^\top\Sigma\beta$. Completing the square
with $\Sigma\Sigma^\dagger c=c$ gives \eqref{eq:projection}.
\end{proof}
The same identity holds conditionally on $\mathcal T$ when
$\E[Q\mid\mathcal T]=0$ and the relevant conditional second moments are
finite, using the conditional payoff mean and conditional moments in
place of $p_N,\Sigma,c$. It is a population identity. Finite-sample
ridge coefficients need not attain it. No retained-variance perturbation
bound relative to a finite-signature benchmark is assumed here.

\subsection{Partial-information quadratic hedging}
Let $\widetilde S_k\in\R^a$ be discounted prices of the $a$ traded risky
assets. To describe the Heston experiments, write $v_k$ for the latent
variance vector and $O_k=v_k+\eta\xi_k$ for its noisy signal at positive
grid times, where $\xi_k$ are independent standard Gaussian observation
errors. The known initial signal equals the prescribed initial variance.
Set
\[
 \mathcal G_k=\sigma(\widetilde S_j,O_j:0\leq j\leq k),\qquad
 \mathcal F_k^{\rm state}=\mathcal G_k\vee\sigma(v_j:0\leq j\leq k).
\]
The simulator filtration additionally contains the shocks used through $k$;
thus $\mathcal G_k\subseteq\mathcal F_k^{\rm state}
\subseteq\mathcal F_k^{\rm sim}$. Completed filtrations can be used without
changing the statements. Let $H\in L^2(\mathcal G_N,\mathbb Q)$ be the
discounted target payoff. An admissible observed-information strategy
$\vartheta_k\in\R^a$ is $\mathcal G_k$-measurable and satisfies
\[
 \sum_{k=0}^{N-1}\E\bigl[(\vartheta_k\cdot
 (\widetilde S_{k+1}-\widetilde S_k))^2\bigr]<\infty.
\]
Its terminal wealth and risk are
\begin{equation}\label{eq:hedging}
 W_N(c,\vartheta)=c+\sum_{k=0}^{N-1}\vartheta_k\cdot
 \Delta\widetilde S_{k+1},\qquad
 \mathcal R(c,\vartheta)=\E[(H-W_N(c,\vartheta))^2].
\end{equation}
Here $c\in\R$ is deterministic initial capital. If a simulation law or
payoff lacks these second moments, the following $L^2$ statements require
an explicit integrability restriction.

For the risk-neutral log-price Euler step used in the experiment, each
discounted asset has conditional multiplicative increment
$\exp(-v_k\Delta t/2+\sqrt{v_k\Delta t}\,\zeta_{k+1})$, with the
nonnegative variance used by the scheme and a standard normal next price
shock conditional on simulator history. Its conditional expectation is one.
Under integrability the discounted price is a simulator martingale and,
by the tower property, a martingale in both smaller filtrations. Correlation
with the next variance shock does not change the conditional Gaussian
marginal of this price shock.

\begin{proposition}[Nested population information spaces]\label{prop:hedging}
Let $H\in L^2(\mathcal G_N,\mathbb Q)$, and suppose $\widetilde S$ is an
adapted square-integrable martingale in the simulator filtration and is
adapted to each of the nested filtrations above. For each such filtration
$\mathcal H$, a strategy is admissible when every $\vartheta_k$ is
$\mathcal H_k$-measurable and
$\sum_{k=0}^{N-1}\E[(\vartheta_k\cdot\Delta\widetilde S_{k+1})^2]<\infty$.
Form the closure $\mathcal K_{\mathcal H}$ in $L^2(\mathbb Q)$ of
constants plus terminal gains of these strategies. Let
$P_{\mathcal H}$ denote orthogonal projection onto this closed subspace.
Then the infimum of \eqref{eq:hedging} for information $\mathcal H$ is
$\norm{H-P_{\mathcal H}H}_2^2$. If
$\mathcal K_1\subseteq\mathcal K_2$, then
\begin{equation}\label{eq:nesting}
 \norm{H-P_1H}_2^2-\norm{H-P_2H}_2^2
 =\norm{P_2H-P_1H}_2^2\geq0.
\end{equation}
In particular the observed-information population risk cannot be below the
state-information population risk. The same statement applies to a closed
linear space of gains generated by fixed observed features contained in
$\mathcal K_{\mathcal G}$.
\end{proposition}
\begin{proof}
Predictable strategies for a smaller filtration are predictable for a larger
one, with the same traded gains and integrability criterion. This proves
inclusion before and after closure. For completeness, projection onto a
closed subspace exists because the parallelogram identity makes a
distance-minimizing sequence Cauchy; its limit lies in the subspace.
Differentiating the squared distance along each subspace direction makes
the residual orthogonal to that subspace and proves uniqueness. For nested
subspaces, $P_1H=P_1P_2H$. Decompose
$H-P_1H=(H-P_2H)+(P_2H-P_1H)$ into orthogonal terms to obtain
\eqref{eq:nesting}.
\end{proof}
Closure makes the infimum precise without asserting that it is attained by
a particular finite readout or predictable integrand. This is the projection
setting underlying martingale quadratic hedging and GKW decompositions.
Here GKW abbreviates Galtchouk--Kunita--Watanabe.
Restricted-information formulations are developed in
\cite{Schweizer,Mania,FreyRunggaldier}; their trading and observation models
are not all the fixed-grid model used here. Simulator innovations in
\eqref{eq:control} are not added to the traded-gain spaces in
\eqref{eq:nesting}. In particular hidden Brownian shocks are not available
to an admissible observed-information hedge.

\subsection{Preliminary experiments and the resulting question}
The implemented reservoir processes a causally available path $U_k$ using
fixed random matrices and biases:
\begin{equation}\label{eq:discrete}
 R_{k+1}=R_k+\sum_{i=1}^d\tanh(A_iR_k+b_i)\Delta U_{k+1}^i,
 \qquad R_0=0.
\end{equation}
The activation acts componentwise. Readouts use time and reservoir state;
hedging features also retain the current observed quantities. Reservoir and
preprocessing choices are fixed before independent test evaluation. This
discrete recurrence motivates the continuous architecture studied in this paper,
but no theorem transferring identifiability to \eqref{eq:discrete} is proved.

For an independent test sample, let $\widehat{\mathcal R}$ be the sample
mean of squared terminal hedging errors in \eqref{eq:hedging}. Its square
root is the reported root mean squared error (RMSE); the sample mean
itself is the mean squared error (MSE).

\begin{table}[htbp]
\centering
\caption{Preliminary independent-test hedging RMSE, quoted from the June
2026 experiment note \cite{JuneNote}. Noise $\eta=0.02$; 10,000/3,000/10,000
training/validation/test paths; validation-only model selection; $N=252$;
one seed. The particle filter was available only in the one-asset case.}
\label{tab:finance}
\begin{tabular}{@{}lrr@{}}
\toprule
Features/model & One-asset Asian & Three-asset basket Asian \\
\midrule
Raw & 2.4037 & 2.0981 \\
Reservoir & 1.0167 & 1.0573 \\
Finite signature & 1.1026 & 1.2486 \\
Particle filter, quadratic readout & 2.2579 & unavailable \\
Causal Transformer & 0.8102 & 0.8684 \\
Fitted true-state model & 0.7921 & 0.8496 \\
\bottomrule
\end{tabular}
\end{table}

The values in \cref{tab:finance} are quoted from the June 2026 experiment
note. Finance code and a text extraction of the note are available locally,
but no run-specific data or fitted-model artifact establishing these exact
measurements was located. They were not regenerated for this paper.
The note reports
similar ordering at $N=52$ and descriptive reservoir gap closures of
$92.1\%$ and $89.2\%$ in the two $N=252$ problems. Gap closure is defined,
when its denominator is positive, by
\[
 \frac{\widehat{\mathcal R}_{\rm raw}-\widehat{\mathcal R}_{\rm reservoir}}
 {\widehat{\mathcal R}_{\rm raw}-\widehat{\mathcal R}_{\rm fitted\ state}}.
\]
The risks here are empirical MSEs, not RMSE differences. Recalculation from
the displayed RMSEs agrees with both percentages to the stated precision;
this checks arithmetic, not experimental reproduction. These percentages
are descriptive and are not an instance of \eqref{eq:nesting}: the
validation-selected fitted true-state model is neither the population
state-information optimum nor assumed to belong to a nested fitted class.
Missing uncertainty intervals are not supplied retrospectively.

The note also reports qualitative variance reduction in preliminary pricing
experiments. Both sets of experiments remain single-seed evidence requiring
further replication before a benchmark conclusion. In the displayed hedging
results the reservoir improves on raw features and finite signatures but has
larger RMSE than the causal Transformer. Improvement over raw features
suggested useful path-history information. It prompted the intrinsic
representation question; the mathematical results in this paper do not explain
these empirical inequalities.

\section{Related work and the contribution boundary}
\label{sec:related}
\subsection{Classical algebra, signatures, and reconstruction}
Free-Lie tensor embedding, Hall coordinates, and the Witt dimensions are
classical \cite{Reutenauer}. The truncated tensor system is also classical
\cite{FrizVictoir,AkyildirimAnomaly}. Finite-dimensional actions of
nilpotent Lie algebras can be constructed from weighted quotients of
universal enveloping algebras; Burde, Eick, and de Graaf review this route
in their Section 2, Theorem 2.1 \cite{Burde}. Their action is expressed
using left multiplication and the usual matrix commutator. Our
right-concatenation orientation is checked directly in
\eqref{eq:linearwords}, including the reversal of matrix composition
caused by the vector-field convention. The existence of the linear
module is not the shallow-tanh transfer argument.

The grading-derived affine action in \cref{lem:gradingaffine} is
classical pre-Lie machinery. Burde's nonsingular-cocycle construction
and its invertible-derivation and positive-grading corollaries
\cite[Section 3.1, Proposition 3.3, Corollaries 3.4--3.5]{BurdePreLie}
supply $\mathcal E^{-1}\operatorname{ad}_a\mathcal E$.
The finite-jet realization by shallow tanh and the resulting fixed-state
Hall minor are proved separately in \cref{sec:hallrank}.

Qualitative rank existence also follows by rescaling this affine action.
At width $\nu$, take $c=1/2$ and
$V_i^\epsilon(x)=\tanh(\epsilon ce_i-\Lambda_{e_i}x)$.
Under $x=\epsilon y$, the pulled-back fields are
$\epsilon^{-1}\tanh(\epsilon W_{e_i}^c(y))$, which converge to
$W_{e_i}^c$ in every fixed finite jet on compact neighborhoods.
Naturality of brackets gives a Hall matrix at zero of the form
$\epsilon(cI_\nu+o(1))$; thus sufficiently small
finite $\epsilon>0$ already gives full rank. This is a direct scaling
observation, not an identified prior publication. The explicit witness
in \cref{lem:hallwitness} additionally fixes $q$, preserves the graded
blocks exactly and supplies the exact selected minors.

Cuchiero, Larsson, and Teichmann
\cite[Theorem 5.1 and Lemma 6.1]{CuchieroLarssonTeichmann} consider
state dimension at least two and at least five vector fields on a bounded
open connected domain. For jointly analytic families with a parameter law
having a density and a specialization whose Lie algebra contains all
polynomial vector fields, they prove interpolation of generic training
inputs with distinct targets, using trained controls and a rescaled
displacement readout. Their proof constructs nonzero
minors of iterated-bracket evaluations and uses analytic zero sets.
Corollary 5.4 and Example 5.5 give seven fields with randomized analytic activations
with linear and quadratic specializations, including mixtures with tanh.
They do not give a prescribed-degree maximal Hall-rank law for fixed
componentwise tanh. The analytic-minor method and the use of a witness
are established tools; our task is to construct the particular finite
tanh minor.

Finite evaluation and local inverse methods also predate this work.
Bailleul and Diehl's Theorem 6 \cite[p.~7]{BailleulDiehl} uses a
reconstruction matrix of generator and degree-two bracket evaluations
at a collection of initial states, including the one-state case.
Their conclusion concerns small-time rough-flow reconstruction in a
specified roughness regime. It does not supply an all-$L$ shallow-tanh
rank witness. Our main observation is the degree-$L$ truncated
polynomial at zero; the separate scalar approximate result concerns
eight actual BV-driven CDE responses. These observation models are
not identified with their rough-flow inverse problem.

\subsection{Randomized signatures and the shallow architecture}
Monzio Compagnoni et al.\ \cite[Definition II.5]{MonzioCompagnoni} use
the same fields $\sigma(A_i x+b_i)$. Their Theorem IV.1 (PDF p.~4)
states a density result in state dimension at least two for
Brownian-driven responses, for almost every initial state, under analytic
activation and i.i.d.\ weights and shifts with absolutely continuous laws; its
proof asserts generic H\"ormander bracket generation. This is a direct
rank antecedent. Unrestricted bracket generation does not determine
the rank at each prescribed degree cutoff or give the minor used here.
We do not use the theorem in its full stated generality as a premise:
analyticity alone allows constant activations, which cannot span a
state space of dimension greater than one. This qualification supplies
no evidence of novelty for the present result.

Grong and Vega-Molino \cite[Remark 3.2]{GrongVegaMolino} use
$\partial_x$ and $\tanh(x)\partial_x$, generate polynomial fields in
$\tanh x$ by brackets, and use Vandermonde evaluations to obtain
controllability of any finite ordered scalar landmark configuration.
Thus both the scalar pair and its polynomial reduction have a precedent.
Their unrestricted Lie spanning does not specify the homogeneous
degree-four Hall maps, their rank-seven mixed-degree relation, or the
eight state-scale decoder studied here. Those are the additional scalar
statements proved in \cref{sec:scalar,sec:cross,sec:eight,sec:stability}.

Cuchiero, Gonon, Grigoryeva, Ortega, and Teichmann develop discrete-time
signature state-affine systems and their random projections as reservoir
approximators \cite{Cuchiero}. Their approximation results for a different
discrete system do not establish the fixed finite word determinant of
\cref{thm:width}. Akyildirim, Gambara, Teichmann, and Zhou use randomized
signature features in market anomaly detection and portfolio selection
\cite{AkyildirimAnomaly,AkyildirimPortfolio}. These works provide
architectural and application antecedents.

Biagini, Gonon, and Walter \cite{BiaginiGononWalter} use a structured
discrete-time randomized-signature construction. Under Sampling
Scheme 4 and Assumption 6, their Lemma 7 gives an almost-sure
continuous injective encoding of a finite datastream in a designated
memory block. Their Proposition 9 gives uniform approximation of continuous
functions of the preterminal history $(x_1,\ldots,x_{T-1})$ by linear
readouts of the terminal increment at sufficient width and prescribed
confidence; it does not allow unrestricted dependence on $x_T$.
These locators refer to the cited PDF version.
The activation assumptions include tanh, but the structured recurrence
and sampling family differ from the full biased-tanh CDE parameter
space in \cref{thm:hallrank}. Their injective encoding does not supply
our Hall evaluation minor; conversely, our local nonlinear inverse
does not establish their linear-readout approximation conclusion.

The market-anomaly paper records the linear tensor system in Proposition
2.10 and states a broad reconstruction claim in Theorem 2.13. Its
analytic-genericity sketch does not supply the specific nonzero witness
needed here. Its activation assumptions also do not directly give a
tanh theorem: the complex Taylor series of tanh does not have infinite
radius of convergence. We therefore use neither that broad claim nor
an all-degree interpretation as a premise of this paper.

Gl\"uckstad, Mu\c{c}a Cirone, and Teichmann \cite{Gluckstad} analyze
iterated-field independence through tree expansions. Their Lemma 3.4 and
Proposition 3.6 give a conditional route from per-degree independence of
complete tree fields to per-degree independence of iterated word operators.
Section 3.3 presents a depth-two nested-exponential construction and claims
that its tree condition holds at state dimension at least $L-1$, with
positive state dimension understood. The intermediate independence
assertion for direction-resolved projected tree terms in version 1 requires
qualification. Dependence among those individual projected terms does not,
by itself, decide independence of the complete summed tree fields. We
therefore do not treat the claimed depth-two width as an independently
verified comparison bound or use that argument as a premise.

Remark 3.10 concerns direction-resolved tree terms for shallow bias-free
fields. Dependence at that level does not imply dependence of the
word-response fields studied here. Section 5 discusses direct
Lie-coordinate recovery. The proof of their Theorem 4.1 uses scale
derivatives and multiple initial states, although scale access is not
explicit in its statement. We keep access to input scaling explicit
where it is used and prove the coefficient formula independently;
no inference from unscaled flows to scale jets is imported. No ranking of the architectures or
refutation of the complete depth-two reconstruction theorem is claimed.

The remaining contribution is an explicit finite-parameter tanh
realization with exact maximal Hall minors at every prescribed finite
degree and width, and the resulting local inverse of the truncated
response at its dimension threshold. The cited rank and interpolation
results do not supply this specified finite-degree matrix. Fixing the
state alone is not the contribution, and affine actions, analytic
genericity, finite evaluation and local inversion are established methods.
The scalar certificates and quantitative decoder, and the larger-width
word transfer, remain separate results with their own assumptions.
Absence of a matching statement in the sources checked, or a qualification
of another proof, does not establish worldwide priority.

\subsection{Financial antecedents}
Martingale control variates, including stochastic-volatility pricing
applications, are established in work such as Fouque and Han
\cite{FouqueHan}. Restricted-information quadratic hedging is studied by
Mania, Tevzadze, and Toronjadze \cite{Mania}. Schweizer \cite{Schweizer}
develops restricted-information risk minimization using GKW and projection
methods. Frey and Runggaldier \cite{FreyRunggaldier} treat stochastic
volatility observed at discrete random times and the associated filtering
problem. Local risk minimization, general mean-variance hedging, and the
fixed-grid martingale projection in \cref{sec:finance} should not be
identified without checking their assumptions. These references establish
the financial background; neither the pricing nor the hedging projection
facts in this paper are presented as a contribution to that theory.

\section{Scope and limitations}
\label{sec:limitations}
The main theorem fixes initialization at zero and concerns Hall-field
evaluations in the full biased parameter space. Its inverse is local
and applies to the degree-$L$ truncated polynomial response.
The sharp width $\nu_{d,L}$ is specific to continuous local recovery
of an open Hall-coordinate neighborhood from one vector.
Neither the rank law nor the local inverse gives global recovery,
a trained linear-readout theorem, or practical conditioning of randomly
sampled parameters. The inverse neighborhood can depend on the
parameters and the nonzero scale.

The three observation objects remain distinct. Ideal scale derivatives
yield exact graded coefficients. For the original scalar witness,
the eight truncated observations have a rational stability certificate,
and \cref{thm:actual} controls the full-CDE remainder and extends its
decoder to perturbed observations by clipping. Its conclusion is
approximate. The stated estimate guarantees consistency as $h$ tends to
zero when the observation error and exact iteration residual are $o(h^4)$.
Exact recovery from actual
responses at a fixed nonzero scale is not established. No corresponding
actual-response bound is asserted for the new Hall-width construction.

The larger width $D_0$ belongs to the all-word independence theorem.
Its all-word field-independence conclusion is not replaced by the Hall-width
result. Neither construction inherits the old scalar radius or inverse
constants. BV signature realization is used for unrestricted paths;
restricted path families, such as a prescribed monotone time channel,
are not covered by the dimensional sharpness statement without their
own open-coordinate-domain argument.

The implemented recurrence \eqref{eq:discrete} is another distinct
object. Its path-feature behavior is not proved by the continuous-field
arguments. In particular, the preliminary hedging improvement is not
a consequence of the Hall-rank theorem.

\section{Conclusion}
Shallow biased tanh attains the maximal fixed-zero Hall evaluation rank
$\min(D,\nu_{d,L})$ generically, with an explicit finite witness.
This gives local analytic recovery of the truncated Hall coordinates
from one zero-initialized vector at the sharp width
$D=\nu_{d,L}$. The specific content is the prescribed-degree tanh
realization, its exact minors and its truncated observation model.
The proof uses classical grading-derived affine actions and the
established analytic-minor method, with a finite-jet argument supplying
the exact tanh blocks. Prior bracket-generation, interpolation and
landmark-control results supply the relevant antecedents
\cite{MonzioCompagnoni,CuchieroLarssonTeichmann,GrongVegaMolino}.
The scalar homogeneous and cross-degree certificates, the ideal
scale-jet decoder, the quantitative eight-observation scalar inverse
and its approximate actual-response extension, and the larger-width
word-independence theorem retain their separate roles and assumptions.

\appendix
\section{Exact scalar graded certificates and reconstruction checks}
\label{app:graded}
The homogeneous matrices and determinants are \eqref{eq:Ms} and
\eqref{eq:dets}. Their columns follow the deterministic Lyndon order in
\eqref{eq:Hall8}; their rows are the stated ordered initial points.
The code uses zero-based generator indices, so internal word $(0,0,1)$
is displayed as $112$, with the same chronological ordering. Exact jets
are normalized Taylor coefficients. The following inverse matrices are
transcribed from the reproduced graded certificate.

\[
M_1^{-1}=\begin{pmatrix}
-2 & 2 \\
2 & 0
\end{pmatrix}.
\]
\[
M_2^{-1}=\begin{pmatrix}
-2
\end{pmatrix}.
\]
\begingroup
\renewcommand{\arraystretch}{1.6}
\[
M_3^{-1}=\begin{pmatrix}
\frac{16}{15} & \frac{16}{15} \\
\frac{8}{3} & -\frac{8}{3}
\end{pmatrix}.
\]
\[
M_4^{-1}=\begin{pmatrix}
-\frac{64}{45} & \frac{8}{15} & -\frac{64}{45} \\
-\frac{16}{3} & 0 & \frac{16}{3} \\
-\frac{128}{45} & \frac{76}{15} & -\frac{128}{45}
\end{pmatrix}.
\]
\endgroup

These inverses implement \eqref{eq:recursive}, with the ordered tensor
corrections in \eqref{eq:exp4}. The first exact check uses
\[
 z=((1,-1),(2),(3,-2),(-1,2,-3)).
\]
Its observations, corrections, and recovered coordinates are shown below.
No signature simulation or floating-point matrix inverse is used in these
checks: each formal tensor exponential and jet evaluation is rational.

\begin{center}
\small
\renewcommand{\arraystretch}{1.7}
\begin{tabular}{@{}rllll@{}}
\toprule
$m$ & $j_m$ & $\psi_m$ & $j_m-\psi_m$ & $z_m$ \\
\midrule
1 & $\bigl(-\frac{1}{2},0\bigr)$ & $\bigl(0,0\bigr)$ & $\bigl(-\frac{1}{2},0\bigr)$ & $\bigl(1,-1\bigr)$ \\
2 & $\bigl(-\frac{5}{4}\bigr)$ & $\bigl(-\frac{1}{4}\bigr)$ & $\bigl(-1\bigr)$ & $\bigl(2\bigr)$ \\
3 & $\bigl(\frac{37}{32},\frac{3}{2}\bigr)$ & $\bigl(\frac{1}{8},-\frac{9}{32}\bigr)$ & $\bigl(\frac{33}{32},\frac{57}{32}\bigr)$ & $\bigl(3,-2\bigr)$ \\
4 & $\bigl(\frac{445}{512},\frac{5}{32},\frac{207}{256}\bigr)$ & $\bigl(\frac{385}{512},\frac{13}{32},\frac{81}{256}\bigr)$ & $\bigl(\frac{15}{128},-\frac{1}{4},\frac{63}{128}\bigr)$ & $\bigl(-1,2,-3\bigr)$ \\
\bottomrule
\end{tabular}
\end{center}

Two further nonzero rational inputs and their observed jets follow. Each
was reconstructed to exactly its input vector. They test the lower-degree
corrections, rather than only the zero-coordinate linearization.

\paragraph{Exact check 2.}
\begin{align*}
z_1&=\bigl(\frac{1}{2},-\frac{2}{3}\bigr),&j_1&=\bigl(-\frac{1}{3},-\frac{1}{12}\bigr),\\
z_2&=\bigl(\frac{5}{7}\bigr),&j_2&=\bigl(-\frac{37}{84}\bigr),\\
z_3&=\bigl(-\frac{3}{5},\frac{4}{9}\bigr),&j_3&=\bigl(-\frac{5239}{32256},-\frac{13817}{32256}\bigr),\\
z_4&=\bigl(\frac{7}{11},-\frac{8}{13},\frac{9}{14}\bigr),&j_4&=\bigl(-\frac{1599438283}{15498362880},-\frac{393733}{1995840},-\frac{6666901981}{15498362880}\bigr).
\end{align*}
\paragraph{Exact check 3.}
\begin{align*}
z_1&=\bigl(-\frac{5}{3},\frac{2}{7}\bigr),&j_1&=\bigl(\frac{1}{7},-\frac{29}{42}\bigr),\\
z_2&=\bigl(-\frac{11}{4}\bigr),&j_2&=\bigl(\frac{211}{168}\bigr),\\
z_3&=\bigl(\frac{13}{8},-\frac{17}{9}\bigr),&j_3&=\bigl(\frac{27269}{84672},\frac{63241}{84672}\bigr),\\
z_4&=\bigl(\frac{19}{10},-\frac{23}{12},\frac{29}{15}\bigr),&j_4&=\bigl(-\frac{74384759}{142248960},-\frac{2116939}{4445280},-\frac{45565913}{17781120}\bigr).
\end{align*}

The cross-degree polynomial recurrence \eqref{eq:polybracket} gives
\eqref{eq:P8} and the complete coefficient matrix \eqref{eq:coeffmatrix}.
As an independent exact convention check, their Taylor expansions after
$u=\tanh x$ through order six agree with the existing recursive bracket-jet
implementation. Because all displayed polynomials have degree at most six
and $\tanh$ has nonzero linear coefficient, this comparison is consistent
with the exact polynomial rank proof. The rank itself is certified by the
nonzero rational minor and explicit kernel vector, not by floating-point
Taylor evaluation.

\section{The complete state-scale Jacobian and inverse}
\label{app:inverse}
With design rows in \eqref{eq:design} and Hall columns in
\eqref{eq:Hall8}, the full derivative of the nonlinear truncated map is
the following rational matrix.

\begingroup\small\setlength{\arraycolsep}{3pt}
\renewcommand{\arraystretch}{1.7}
\begin{equation}\label{eq:Bfull}
B=\begin{pmatrix}
-\frac{1}{2} & \frac{1}{2} & -\frac{3}{8} & \frac{15}{32} & \frac{3}{16} & -\frac{57}{128} & -\frac{3}{32} & \frac{3}{64} \\
-1 & 1 & -\frac{3}{2} & \frac{15}{4} & \frac{3}{2} & -\frac{57}{8} & -\frac{3}{2} & \frac{3}{4} \\
-\frac{3}{2} & \frac{3}{2} & -\frac{27}{8} & \frac{405}{32} & \frac{81}{16} & -\frac{4617}{128} & -\frac{243}{32} & \frac{243}{64} \\
0 & 1 & -2 & 4 & 0 & -8 & 0 & 4 \\
0 & \frac{3}{2} & -\frac{9}{2} & \frac{27}{2} & 0 & -\frac{81}{2} & 0 & \frac{81}{4} \\
\frac{1}{2} & \frac{1}{2} & -\frac{3}{8} & \frac{15}{32} & -\frac{3}{16} & -\frac{57}{128} & \frac{3}{32} & \frac{3}{64} \\
1 & 1 & -\frac{3}{2} & \frac{15}{4} & -\frac{3}{2} & -\frac{57}{8} & \frac{3}{2} & \frac{3}{4} \\
\frac{3}{2} & \frac{3}{2} & -\frac{27}{8} & \frac{405}{32} & -\frac{81}{16} & -\frac{4617}{128} & \frac{243}{32} & \frac{243}{64}
\end{pmatrix}.
\end{equation}
\endgroup

For an additional elimination certificate, ordinary exact LU elimination
swaps the second and fourth rows and has diagonal pivots
\[
 -\frac12,\quad 1,\quad-\frac94,\quad-\frac{15}{16},\quad-3,
 \quad\frac{4863}{1280},\quad-\frac{9339}{3242},
 \quad-\frac{32265}{2264}.
\]
The negative of their product is $-258733035/524288$, agreeing with
\eqref{eq:detB}. The row-scaled integer display below specifies the entire
inverse without rounding: if $\mathbf s$ is its positive denominator
vector and $T_B$ the displayed integer matrix, then
$B^{-1}=\diag(1/s_1,\ldots,1/s_8)T_B$. Direct multiplication on either
side gives the identity.

\[
\mathbf s=\bigl(66,717,2151,2151,99,32265,33,32265\bigr).
\]
\begingroup\small\setlength{\arraycolsep}{3pt}
\[
T_B=\begin{pmatrix}
-108 & 27 & -4 & 0 & 0 & 108 & -27 & 4 \\
1080 & 378 & -88 & -1215 & 240 & 1080 & 378 & -88 \\
-684 & 4923 & -964 & -8910 & 1760 & -684 & 4923 & -964 \\
-2280 & 3504 & -664 & -3888 & 768 & -2280 & 3504 & -664 \\
-228 & 156 & -28 & 0 & 0 & 228 & -156 & 28 \\
-9168 & 12852 & -2992 & -11196 & 2424 & -9168 & 12852 & -2992 \\
-40 & 32 & -8 & 0 & 0 & 40 & -32 & 8 \\
-1416 & 5814 & -2264 & -9162 & 3828 & -1416 & 5814 & -2264
\end{pmatrix}.
\]
\endgroup

In the same row order its exact absolute row sums are
\begin{equation}\label{eq:invrows}
 \left(\frac{139}{33},\frac{4547}{717},\frac{23812}{2151},
 \frac{17552}{2151},\frac{824}{99},\frac{63644}{32265},
 \frac{160}{33},\frac{31978}{32265}\right).
\end{equation}
Their maximum is $K=23812/2151$. This exact infinity norm is distinct
from a numerical two-norm condition number.

\section{Complete polynomial specification and stability arithmetic}
\label{app:polynomial}
The coefficient calculation can be stated without any opaque polynomial
oracle. Let $\mathbf d=(1,1,2,3,3,4,4,4)$, and for a multiindex $\alpha$
put $\mathrm{wt}(\alpha)=\sum_a d_a\alpha_a$. For an initial point $y$,
write
\begin{equation}\label{eq:polyfull}
 F_y^{\leq4}(r;z)-y
 =\sum_{\substack{1\leq|\alpha|\leq4\\\mathrm{wt}(\alpha)\leq4}}
 c_{y,\alpha}\,r^{\mathrm{wt}(\alpha)}z^\alpha.
\end{equation}
Each coefficient has the finite expression
\begin{equation}\label{eq:coeffalgorithm}
 c_{y,\alpha}=\frac1{|\alpha|!}
 \sum_{\substack{(a_1,\ldots,a_{|\alpha|})\\
 \#\{j:a_j=b\}=\alpha_b\ \text{for every }b}}
 \bigl(\Gamma(H_{a_1})\cdots\Gamma(H_{a_{|\alpha|}})\Id\bigr)(y).
\end{equation}
This sums ordered tuples, preserving noncommuting operators even though
the coordinate monomials commute. By \eqref{eq:Lieoperator}, each
$\Gamma(H_a)$ is the scalar operator
$P_a(u)(1-u^2)\,d/du$ on functions of $u$. When the rightmost operator
acts on $\Id(x)=x$, it gives $P_a(u)$; all remaining actions use that
displayed polynomial derivative. Hence \eqref{eq:P8} and
\eqref{eq:coeffalgorithm} determine every entry by rational operations
at $u=-1/2,0,1/2$.

\Cref{tab:responsecoeff} lists all 27 monomials with a nonzero coefficient
at at least one of these states. Unlisted monomials have zero coefficient
at all three states. Substitute each design scale in \eqref{eq:polyfull}
to recover the eight complete observation polynomials; their total number
of nonzero output terms is 190. Differentiating a term contributes
$\alpha_a c_{y,\alpha}r^{\mathrm{wt}(\alpha)}z^{\alpha-e_a}$ to column
$a$ of the Jacobian, where $e_a$ is the coordinate multiindex. This
specifies all coefficients in \eqref{eq:derivative_poly} and directly
connects the majorant to the implemented nonlinear map.

\begingroup
\renewcommand{\arraystretch}{1.7}
\begin{longtable}{@{}lrrrr@{}}
\caption{Complete nonzero coefficient support in \eqref{eq:polyfull}. The scale power is the weight.}\label{tab:responsecoeff}\\
\toprule
$z^\alpha$ & $\mathrm{wt}(\alpha)$ & $c_{-q,\alpha}$ & $c_{0,\alpha}$ & $c_{q,\alpha}$ \\
\midrule
\endfirsthead
\toprule
$z^\alpha$ & $\mathrm{wt}(\alpha)$ & $c_{-q,\alpha}$ & $c_{0,\alpha}$ & $c_{q,\alpha}$ \\
\midrule
\endhead
$z_{2}$ & 1 & $\frac{1}{2}$ & $\frac{1}{2}$ & $\frac{1}{2}$ \\
$z_{1}$ & 1 & $-\frac{1}{2}$ & $0$ & $\frac{1}{2}$ \\
$z_{3}$ & 2 & $-\frac{3}{8}$ & $-\frac{1}{2}$ & $-\frac{3}{8}$ \\
$z_{1}z_{2}$ & 2 & $\frac{3}{16}$ & $\frac{1}{4}$ & $\frac{3}{16}$ \\
$z_{1}^{2}$ & 2 & $-\frac{3}{16}$ & $0$ & $\frac{3}{16}$ \\
$z_{5}$ & 3 & $\frac{3}{16}$ & $0$ & $-\frac{3}{16}$ \\
$z_{4}$ & 3 & $\frac{15}{32}$ & $\frac{1}{2}$ & $\frac{15}{32}$ \\
$z_{2}z_{3}$ & 3 & $-\frac{3}{32}$ & $0$ & $\frac{3}{32}$ \\
$z_{1}z_{3}$ & 3 & $-\frac{3}{64}$ & $-\frac{1}{4}$ & $-\frac{3}{64}$ \\
$z_{1}z_{2}^{2}$ & 3 & $\frac{1}{32}$ & $0$ & $-\frac{1}{32}$ \\
$z_{1}^{2}z_{2}$ & 3 & $-\frac{1}{64}$ & $\frac{1}{12}$ & $-\frac{1}{64}$ \\
$z_{1}^{3}$ & 3 & $-\frac{1}{64}$ & $0$ & $\frac{1}{64}$ \\
$z_{8}$ & 4 & $\frac{3}{64}$ & $\frac{1}{4}$ & $\frac{3}{64}$ \\
$z_{7}$ & 4 & $-\frac{3}{32}$ & $0$ & $\frac{3}{32}$ \\
$z_{6}$ & 4 & $-\frac{57}{128}$ & $-\frac{1}{2}$ & $-\frac{57}{128}$ \\
$z_{3}^{2}$ & 4 & $\frac{9}{128}$ & $0$ & $-\frac{9}{128}$ \\
$z_{2}z_{5}$ & 4 & $-\frac{3}{128}$ & $-\frac{1}{8}$ & $-\frac{3}{128}$ \\
$z_{2}z_{4}$ & 4 & $\frac{3}{64}$ & $0$ & $-\frac{3}{64}$ \\
$z_{1}z_{5}$ & 4 & $\frac{3}{32}$ & $0$ & $-\frac{3}{32}$ \\
$z_{1}z_{4}$ & 4 & $\frac{33}{256}$ & $\frac{1}{4}$ & $\frac{33}{256}$ \\
$z_{2}^{2}z_{3}$ & 4 & $\frac{1}{128}$ & $\frac{1}{24}$ & $\frac{1}{128}$ \\
$z_{1}z_{2}z_{3}$ & 4 & $-\frac{7}{64}$ & $0$ & $\frac{7}{64}$ \\
$z_{1}^{2}z_{3}$ & 4 & $\frac{17}{256}$ & $-\frac{1}{12}$ & $\frac{17}{256}$ \\
$z_{1}z_{2}^{3}$ & 4 & $-\frac{1}{512}$ & $-\frac{1}{96}$ & $-\frac{1}{512}$ \\
$z_{1}^{2}z_{2}^{2}$ & 4 & $\frac{15}{512}$ & $0$ & $-\frac{15}{512}$ \\
$z_{1}^{3}z_{2}$ & 4 & $-\frac{45}{1024}$ & $\frac{1}{48}$ & $-\frac{45}{1024}$ \\
$z_{1}^{4}$ & 4 & $\frac{17}{1024}$ & $0$ & $-\frac{17}{1024}$ \\
\bottomrule
\end{longtable}
\endgroup

\begin{table}[htbp]
\centering
\caption{Exact absolute row coefficient sums in \eqref{eq:Delta}.}\label{tab:majorant}
\renewcommand{\arraystretch}{1.7}
\begin{tabular}{@{}lrrr@{}}
\toprule
Row $j$ & $a_{j,1}$ & $a_{j,2}$ & $a_{j,3}$ \\
\midrule
1,6 & $\frac{225}{128}$ & $\frac{189}{256}$ & $\frac{47}{128}$ \\
2,7 & $\frac{135}{8}$ & $\frac{165}{16}$ & $\frac{47}{8}$ \\
3,8 & $\frac{9369}{128}$ & $\frac{12717}{256}$ & $\frac{3807}{128}$ \\
4 & $18$ & $8$ & $2$ \\
5 & $\frac{315}{4}$ & $\frac{297}{8}$ & $\frac{81}{8}$ \\
\bottomrule
\end{tabular}
\end{table}

The sixteen radii below use the same coefficient-sum bound, with no
floating-point acceptance threshold. The rule selects the largest listed
radius for which the final column is at most $1/2$. Every entry is an
exact rational value; the pass/fail flag is a rational comparison.

\begingroup\small
\renewcommand{\arraystretch}{1.7}
\begin{longtable}{@{}lrrl@{}}
\caption{Complete prescribed exact radius search.}\label{tab:radii}\\
\toprule
$R$ & $\Delta(R)$ & $K\Delta(R)$ & $K\Delta(R)\leq1/2$ \\
\midrule
\endfirsthead
\toprule
$R$ & $\Delta(R)$ & $K\Delta(R)$ & $K\Delta(R)\leq1/2$ \\
\midrule
\endhead
$\frac{1}{2}$ & $\frac{3375}{64}$ & $\frac{2232375}{3824}$ & no \\
$\frac{1}{4}$ & $\frac{11349}{512}$ & $\frac{7506733}{30592}$ & no \\
$\frac{1}{8}$ & $\frac{42777}{4096}$ & $\frac{28294609}{244736}$ & no \\
$\frac{1}{16}$ & $\frac{166113}{32768}$ & $\frac{109874521}{1957888}$ & no \\
$\frac{1}{32}$ & $\frac{654705}{262144}$ & $\frac{433050985}{15663104}$ & no \\
$\frac{1}{64}$ & $\frac{2599569}{2097152}$ & $\frac{1719470473}{125304832}$ & no \\
$\frac{1}{128}$ & $\frac{10360017}{16777216}$ & $\frac{6852575689}{1002438656}$ & no \\
$\frac{1}{256}$ & $\frac{41363793}{134217728}$ & $\frac{27359851081}{8019509248}$ & no \\
$\frac{1}{512}$ & $\frac{165302865}{1073741824}$ & $\frac{109338661705}{64156073984}$ & no \\
$\frac{1}{1024}$ & $\frac{660907089}{8589934592}$ & $\frac{437153322313}{513248591872}$ & no \\
$\frac{1}{2048}$ & $\frac{2643019857}{68719476736}$ & $\frac{1748210800969}{4105988734976}$ & yes \\
$\frac{1}{4096}$ & $\frac{10570862673}{549755813888}$ & $\frac{6992038388041}{32847909879808}$ & yes \\
$\frac{1}{8192}$ & $\frac{42281017425}{4398046511104}$ & $\frac{27966544081225}{262783279038464}$ & yes \\
$\frac{1}{16384}$ & $\frac{169119203409}{35184372088832}$ & $\frac{111862957543753}{2102266232307712}$ & yes \\
$\frac{1}{32768}$ & $\frac{676467081297}{281474976710656}$ & $\frac{447445392773449}{16818129858461696}$ & yes \\
$\frac{1}{65536}$ & $\frac{2705848860753}{2251799813685248}$ & $\frac{1789768696451401}{134545038867693568}$ & yes \\
\bottomrule
\end{longtable}
\endgroup

The value of $c_*$ is exactly $1/K-\Delta(R_*)$ and $C_*=1/c_*$.
The half-margin and the simple bound $2K$ in \cref{thm:stability} are
therefore consequences of the displayed arithmetic. They do not depend
on numerical estimates of a smallest singular value.

\section{Computational verification and reproducibility}
\label{app:reproduce}
\subsection{Authority and exact arithmetic}
The paper's exact scalar data come from the following artifacts:
\begin{itemize}[nosep]
\item \file{research_output/certificate.json}
\item \file{research_output/graded_reconstruction_certificate.json}
\item \file{research_output/finite_observation_certificate.json}
\item \file{research_output/local_stability_certificate.json}
\end{itemize}
Their default reproducers recompute and compare the certificate payloads.
They do not write new artifacts by default. The mathematical results
are proved in \cref{sec:hallrank,sec:graded,sec:finite,sec:width}; a finite computational
test is not substituted for any universal quantifier in those proofs.

The implementation uses the existing deterministic free-Lie words and
Lyndon brackets, exact activation jets, tensor products, and rational
matrix operations. No second free-Lie implementation or tensor convention
is introduced for this paper. The scalar certificate status is
\file{exact_nonzero}, with method
\file{exact_symbolic_determinant_simplification}. The original numerical
candidate remains \file{numerical_candidate_only}.

The separate tensor-system structural harness has the following scope.
\begin{center}
\begin{tabular}{@{}rrrrr@{}}
\toprule
$d$ & $L$ & $D_0$ & $Q_{d,L}$ & Exact Hall rank \\
\midrule
2&2&7&6&3\\
2&3&15&14&5\\
2&4&31&30&8\\
3&2&13&12&6\\
\bottomrule
\end{tabular}
\end{center}
For these cases it verifies the exact identity word matrix, vanishing at
length $L+1$ in the linear system, the one-segment tensor exponential,
the two-segment chronological Chen product, and agreement of Hall
expansions with word responses. Smaller cases also check block embeddings
and report a normalized-tanh limit diagnostic. It extracts exact symbolic
coefficients of orders $\epsilon^0$ and $\epsilon^1$ and represents them
as exact rational values. It certifies no particular nonzero epsilon.
The analytic transfer proof remains the basis for the arbitrary-$d,L$
conclusion.

\subsection{Exact checks of the Hall-rank construction}
The authoritative proof note is
\file{research_output/intrinsic_hall_rank_theorem.md}.
The exact artifact
\file{research_output/intrinsic_hall_rank_certificate.json}
is recomputed from actual finite tanh jets by
\file{research_output/reproduce_intrinsic_hall_rank.py}.
The computations reuse the existing Hall basis and tensor convention.
They construct adjoints, grading operators and ideal quotients exactly;
tanh coefficients are generated recursively from $T'=1-T^2$, with
$\tanh(q)=1/2$ represented exactly. Brackets of length at most $L$
need generator jets through order $L-1$.

For $(d,L,D)=(2,4,8)$, in the order
$(1,2,12,112,122,1112,1122,1222)$, the computed matrix is
\begin{equation}\label{eq:hallregression}
 H_\theta=\frac12 I_8-\frac1{32}(E_{3,6}+E_{3,8}),
 \qquad \det H_\theta=\frac1{256},
\end{equation}
where $E_{a,b}$ has its only nonzero entry, one, in row $a$, column $b$
with one-based indices. This specifies every entry, including the two
nonlinear off-diagonal terms. The calculation generates these entries
from the jets; the displayed matrix is a regression fixture.

The 17 persisted cases include the full-width references
\[
 (d,L,\nu)=(2,2,3),(2,3,5),(2,4,8),(2,5,14),
              (3,3,14),(3,4,32),
\]
as well as widths below $d$, quotients cutting a homogeneous level,
padded widths, and $d=1$ or $L=1$ edge cases. Each selected minor is
exactly $2^{-\min(D,\nu)}$. The finite harness supports
$d\leq3$, $L\leq5$, $d^L\leq128$, and $D\leq40$; these are resource
limits of this implementation, not hypotheses of \cref{thm:hallrank}.
Genericity comes from its proof, not from the finite cases or determinant
sizes. Independent checks include a small complete truncated-map
Jacobian, nonzero rational scale factors and the degenerate zero scale,
direct bracket signs, and the degree-seven tanh coefficient $-17/315$.

\subsection{Non-certifying numerical diagnostics}
The rounded two-norm condition numbers attached to the homogeneous
matrices $M_1,\ldots,M_4$ are respectively
$2.618034$, $1.0$, $2.5$, and $6.2704558$. The certificate's binary64
conversion of the exact determinant product is
$-0.00036209821701049805$. These are diagnostics of exact matrices,
not certificates. The associated numerical ranks are $2,1,2,3$.

The original degree-four numerical discovery candidate used 42 points,
weight seed 1729, point seed 2719, and extension point seed 1000002719.
Its zero-based selected rows were $(18,32,22)$ and columns $(0,1,2)$;
its floating determinant was $-5.041803997872895$ and floating condition
number $1.9857803041316449$. It is a different candidate from the exact
scalar witness. These historical values are preserved as numerical
provenance only.

For the selected eight-output design the binary64 two-norm condition
number reported to twelve significant digits is $468.626825569$.
The reported singular values, in descending order, are
\begin{align*}
 &(72.8678806904,\ 13.4576527068,\ 12.4585453832,\ 2.9110008089,\\
 &\hspace{8mm}1.23631176424,\ 0.365987026873,\ 0.197225996535,
 \ 0.155492337857).
\end{align*}
The deterministic design search first required an exact nonzero
determinant, then preferred fewer distinct absolute scales, a
lexicographically smaller set of absolute scales, a lower numerical
condition number, and finally the lexicographic ordered-pair list.
The numerical tie-break does not enter the determinant or stability proof.

\subsection{Commands and results}
The following commands reproduce the six exact/structural checks:
\begin{verbatim}
python research_output/reproduce_certificate.py
python research_output/reproduce_graded_reconstruction.py
python research_output/reproduce_finite_observation.py
python research_output/reproduce_local_stability.py
python research_output/reproduce_tensor_witness.py
python research_output/reproduce_intrinsic_hall_rank.py
\end{verbatim}
The first five completed successfully in the original manuscript audit.
The Hall reproducer checks all 17 persisted cases without writing by
default. Focused tests are
\begin{verbatim}
python -m pytest tests/test_certificates.py -q
python -m pytest tests/test_graded_reconstruction.py -q
python -m pytest tests/test_finite_observation.py -q
python -m pytest tests/test_local_stability.py -q
python -m pytest tests/test_tensor_witness.py -q
python -m pytest tests/test_intrinsic_hall_rank.py -q
python -m pytest -q
git diff --check
\end{verbatim}
The original 5 September audit recorded respectively 7, 28, 21, 18,
and 40 passing tests for the first five focused commands, and 233
passing tests in 118.74 seconds for its full suite. Those counts are
historical. Actual test
invocations added \texttt{-p no:cacheprovider} and set
\file{PYTHONDONTWRITEBYTECODE=1} to avoid cache artifacts. The environment
used Python 3.13.3, NumPy 2.3.1, SymPy 1.14.0, pytest 9.1.1, and
PyTorch 2.12.1 on Windows.

The completed Hall implementation stage recorded 38 new focused tests,
135 affected existing tests, and a full-suite result of 271 passed.
Its exact commands, timings and environment are retained in the
authoritative theorem note. The Stage 3 document integration reran
the Hall and scalar certificate reproducers and checked document
consistency without repeating the full computational campaign.
The final independent audit reran all six reproducers and the repository
gate \texttt{python -B -m pytest -q -p no:cacheprovider}: all 271 tests
passed in 127.32 seconds. Its additional direct symbolic bracket check
recovered the complete width-eight Hall matrix, and a corrupted payload
with unchanged rank and determinant was rejected by the read-only verifier.

At that audit, the worktree already had unrelated finance and line-ending
changes. Global whitespace checking reported 804 pre-existing
trailing-whitespace instances in \file{rsmc/__init__.py},
\file{rsmc/hedging.py}, and \file{tests/test_hedging.py}.
Those files were not changed during paper preparation. The companion
\file{research_output/paper_mathematical_audit.md} records the source
hierarchy, theorem classifications, and final preservation checks.
The subsequent actual-response validation uses the existing tensor and
polynomial utilities and writes no artifacts. Run
\begin{verbatim}
python research_output/validate_actual_response_recovery.py
\end{verbatim}
It checks exact polynomial scaling and reference Hall conversion, then
compares the axis-flow responses and clipped decoder at 80 and 120 decimal
digits. The existing research report records the paths, perturbations,
remainder, coordinate error and the two distinct residuals. These numerical
flow evaluations and computed residuals are not exact certificates.

Neither a Windows nor an Ubuntu-on-WSL LaTeX engine was available at the
Stage 3 integration check on 13 September; its static checks did not
verify PDF layout. The final independent audit used a portable Tectonic
0.17.0 executable, verified against the official release checksum, to
compile the canonical source. The companion audit records the resulting
diagnostics and visual checks. With Tectonic available on the executable
path, ensure \file{paper/build} exists and run:
\begin{verbatim}
tectonic --keep-logs --keep-intermediates --outdir paper/build \
         paper/randomized_signature_identifiability.tex
\end{verbatim}
Tectonic reruns TeX until the references settle; repeat the command after
any source change.

\begin{thebibliography}{99}

\bibitem{Cuchiero}
C. Cuchiero, L. Gonon, L. Grigoryeva, J.-P. Ortega, and J. Teichmann.
Discrete-time signatures and randomness in reservoir computing.
\emph{IEEE Transactions on Neural Networks and Learning Systems},
33(11):6321--6330, 2022.
\url{https://arxiv.org/abs/2010.14615}.

\bibitem{AkyildirimAnomaly}
E. Akyildirim, M. Gambara, J. Teichmann, and S. Zhou.
Applications of signature methods to market anomaly detection.
Preprint, 2022. \url{https://arxiv.org/abs/2201.02441}.

\bibitem{AkyildirimPortfolio}
E. Akyildirim, M. Gambara, J. Teichmann, and S. Zhou.
Randomized signature methods in optimal portfolio selection.
\emph{Quantitative Finance}, 25(2):197--216, 2025.
\url{https://doi.org/10.1080/14697688.2025.2458613}.

\bibitem{MonzioCompagnoni}
E. Monzio Compagnoni, A. Scampicchio, L. Biggio, A. Orvieto,
T. Hofmann, and J. Teichmann.
On the effectiveness of Randomized Signatures as Reservoir for Learning
Rough Dynamics. \emph{International Joint Conference on Neural Networks
(IJCNN)}, 2023.
Numbered locators refer to arXiv version 3, 26 April 2023:
\url{https://arxiv.org/pdf/2201.00384v3}.

\bibitem{CuchieroLarssonTeichmann}
C. Cuchiero, M. Larsson, and J. Teichmann.
Deep neural networks, generic universal interpolation, and controlled ODEs.
\emph{SIAM Journal on Mathematics of Data Science}, 2(3):901--919, 2020.
\url{https://doi.org/10.1137/19M1284117}.
Numbered locators refer to arXiv version 2:
\url{https://arxiv.org/pdf/1908.07838v2}.

\bibitem{GrongVegaMolino}
E. Grong and S. Vega-Molino.
Controllability on landmark manifolds for shapes and neural ODEs.
\emph{Mathematics of Control, Signals, and Systems}, 38:275--289, 2026.
\url{https://doi.org/10.1007/s00498-026-00435-1}.
Numbered locators refer to arXiv version 3, 31 July 2025:
\url{https://arxiv.org/html/2403.08090v3}.

\bibitem{Gluckstad}
M. Gl\"uckstad, N. Mu\c{c}a Cirone, and J. Teichmann.
Signature reconstruction from randomized signatures.
Preprint, version 1, 2025.
\url{https://arxiv.org/abs/2502.03163v1}.

\bibitem{BailleulDiehl}
I. Bailleul and J. Diehl.
The inverse problem for rough controlled differential equations.
Preprint, version 3, 14 November 2014.
\url{https://arxiv.org/pdf/1407.2768v3}.

\bibitem{BiaginiGononWalter}
F. Biagini, L. Gonon, and N. Walter.
Universal randomised signatures for generative time series modelling.
Preprint, version 2, 6 September 2024.
Numbered locators refer to the PDF:
\url{https://arxiv.org/pdf/2406.10214v2}.

\bibitem{Reutenauer}
C. Reutenauer. \emph{Free Lie Algebras}.
London Mathematical Society Monographs, New Series 7.
Oxford University Press, 1993.
\url{https://academic.oup.com/book/53911}.

\bibitem{FrizVictoir}
P. K. Friz and N. B. Victoir.
\emph{Multidimensional Stochastic Processes as Rough Paths: Theory and
Applications}. Cambridge University Press, 2010.
Numbered locators here refer to the authors' May 2009 draft:
\url{https://page.math.tu-berlin.de/~friz/master4_May6th.pdf}.

\bibitem{Burde}
D. Burde, B. Eick, and W. de Graaf.
Computing faithful representations for nilpotent Lie algebras.
\emph{Journal of Algebra}, 322:602--612, 2009.
\url{https://arxiv.org/abs/0807.2345}.

\bibitem{BurdePreLie}
D. Burde.
Left-symmetric algebras, or pre-Lie algebras in geometry and physics.
\emph{Central European Journal of Mathematics}, 4(3):323--357, 2006.
Numbered locators refer to the author's PDF dated 1 March 2011:
\url{https://homepage.univie.ac.at/dietrich.burde/papers/burde_24_pre_lie.pdf}.

\bibitem{Mityagin}
B. S. Mityagin. The zero set of a real analytic function.
Preprint, 2015. \url{https://arxiv.org/abs/1512.07276}.

\bibitem{FouqueHan}
J.-P. Fouque and C.-H. Han.
A martingale control variate method for option pricing with stochastic
volatility. \emph{ESAIM: Probability and Statistics}, 11:40--54, 2007.
\url{https://www.numdam.org/item/10.1051/ps:2007005.pdf}.

\bibitem{Mania}
M. Mania, R. Tevzadze, and T. Toronjadze.
Mean-variance hedging under partial information.
Preprint, 2007. \url{https://arxiv.org/abs/math/0703424}.

\bibitem{Schweizer}
M. Schweizer. Risk-minimizing hedging strategies under restricted
information. \emph{Mathematical Finance}, 4:327--342, 1994.
\url{https://people.math.ethz.ch/~mschweiz/Files/restrinf.pdf}.

\bibitem{FreyRunggaldier}
R. Frey and W. J. Runggaldier.
Risk-minimizing hedging strategies under restricted information:
The case of stochastic volatility models observable only at discrete
random times. \emph{Mathematical Methods of Operations Research},
50:339--350, 1999.
\url{https://doi.org/10.1007/s001860050101}.

\bibitem{Hatcher}
A. Hatcher. \emph{Algebraic Topology}.
Cambridge University Press, 2002.
\url{https://pi.math.cornell.edu/~hatcher/AT/AT.pdf}.

\bibitem{JuneNote}
P. Stojkovi\'{c}. \emph{Updated Research Note: Randomized-Signature
Reservoirs for Partial-Information Hedging}.
Subtitle: Causal State Representations for Observed-Filtration GKW Hedging.
Internal experiment note, June 2026.
Repository text extraction:
\file{research_output/sources/research_notes_innovation_project_6.txt}.
The displayed financial measurements are quoted from this note.

\end{thebibliography}
\end{document}
